% concatenated from main_arxiv.tex, meta_arxiv.tex, sections/abstract.tex, sections/introduction.tex, sections/model.tex, sections/policies.tex, sections/empirical_analysis.tex, sections/real_world_extension.tex, sections/discussion.tex, sections/acknowledgements.tex, sections/appendix.tex
\documentclass[11pt]{article}
\usepackage[margin=1.2in]{geometry}
\usepackage[noblocks]{authblk}
\usepackage{hyperref}       % hyperlinks
\usepackage{url}            % simple URL typesetting
\usepackage{booktabs}       % professional-quality tables
\usepackage{amsfonts}       % blackboard math symbols
\usepackage{nicefrac}       % compact symbols for 1/2, etc.
\usepackage{makecell}
\usepackage{microtype}      % microtypography
\usepackage{xcolor}         % colors
\usepackage{amsmath}
\usepackage{amsthm}
\usepackage{thmtools}
\usepackage{thm-restate}
% \usepackage{acm-ec-25}
\usepackage{bbm}
\usepackage{natbib}
\usepackage{algorithm}
\usepackage{algorithmic}
\usepackage{subfig}
\usepackage{graphicx}


\usepackage{booktabs} % For formal tables
% \usepackage[ruled]{algorithm2e} % For algorithms
% \renewcommand{\algorithmcfname}{ALGORITHM}

% Choose a citation style by commenting/uncommenting the appropriate line:
%\setcitestyle{acmnumeric}
\setcitestyle{authoryear}
\DeclareMathOperator*{\argmax}{arg\,max}
\DeclareMathOperator*{\argmin}{arg\,min}

\newtheorem{lemma}{Lemma}[section]
\newtheorem{lemma2}{Lemma}[section]
\newtheorem{proposition2}{Proposition}[section]
\newtheorem{theorem2}{Theorem}[section]
\usepackage{amsthm}
\newtheorem{example}{Example}[section] 

%%%%%% Title %%%%%%
\title{Assortment Optimization for Patient-Provider Matching}
\newcommand{\nrcomment}[1]{{\color{red} Naveen: {#1}}}
\newcommand{\hwcomment}[1]{{\color{blue} Holly: {#1}}}

%%%%%% Authors %%%%%%
\author[1]{Naveen Raman}
\author[1]{Holly Wiberg}
\date{}

%%%%%% Affiliations %%%%%%
\affil[1]{Carnegie Mellon University}
\affil[1]{\texttt{\{naveenr, hwiberg\}@cmu.edu}}

\begin{document}

 \maketitle
\newif\ifshowdetails
\showdetailstrue  % or \showdetailsfalse


\begin{abstract}
Rising provider turnover results in frequently needing to rematch patients with available providers. 
However, the rematching process is cumbersome for both patients and health systems, resulting in labor-intensive and ad hoc reassignments. 
We propose a novel patient-provider matching approach to address this issue by offering patients limited provider menus. The goal is to maximize match quality across the system while preserving patient choice. 
We frame this as a novel variant of assortment optimization, where patient-specific provider menus are offered upfront, and patients respond in a random sequence to make their selections. 
This hybrid offline-online setting is understudied in previous literature and captures system dynamics across various domains. 
We first demonstrate that a greedy baseline policy--which offers all providers to all patients--can maximize the match rate but lead to low-quality matches. 
Based on this, we construct a set of policies and demonstrate that the best policy depends on problem specifics, such as a patient's willingness to match and the ratio of patients to providers. 
On real-world data, our proposed policy improves average match quality by 13\% over a greedy solution by tailoring assortments based on patient characteristics. 
Our analysis reveals a tradeoff between menu size and system-wide match quality, highlighting the value of balancing patient choice with centralized planning. 
\end{abstract}


\section{Introduction}
\label{sec:intro}
Primary care providers (PCPs) are essential to the healthcare ecosystem because they are the first point of contact for many patients~\citep{physician_trust,physician_trust_adherence}. 
Patients rely on PCPs for routine checkups and referrals to specialists. PCPs serve as a regular touchpoint for patients and provide care continuity, which has been shown to instill trust and improve medication uptake rates and patient health~\citep{physician_trust_adherence}. Furthermore, high-quality patient-provider matches compound the benefit of primary care, where the notion of a high-quality match is heterogeneous across patients and can include factors such as geographical proximity, provider specialties, and concordance along language, race, and gender~\citep{race_concordance,gender_concordance,language_concordance}. As one such example, \citet{concordance_race_cardio} found that access to race-concordant doctors could reduce cardiovascular disease by 19\%~\citep{concordance_race_cardio}. 

Unfortunately, high provider turnover rates frequently lead patients to lose their PCP, which disrupts patient care and can lead to worse health outcomes~\citep{pcp_turnover}. 
In principle, healthcare administrators reassign unmatched patients to other providers. However, in practice, the process takes months due to provider scarcity and the administrative burden of patient coordination~\citep{finding_new_provider}.  

As an alternative to manual reassignments, algorithmically matching patients and providers can reduce logistical hassle. However, this approach requires carefully balancing patient autonomy and system-wide utility. 
While automatically assigning each patient to a provider would decrease wait times, it also eliminates patient choice~\citep{patient_autonomy,gaynor_free_2016}. 
In contrast, offering patients full autonomy to choose between all available providers could lead to suboptimal matches when considering system-wide match quality; for example, patients who choose quickly might prevent better matches for late-responding patients. In addition, granting patients full independence can be overwhelming and lead to decision delay~\citep{bate_choice_2005}. 

In this work, we formulate the patient-provider matching problem as a variant of assortment optimization~\citep{assortment_school,assortment_dating,assortment_mnl}; we illustrate this problem in Figure~\ref{fig:pull}.
Assortments are menus of options given to customers, who then make selections from this menu~\citep{assortment_optimization}. This framework is commonly seen on e-commerce platforms~\citep{e_commerce_assortment}. 
By offering limited provider menus, patients have the autonomy to select their preferred provider, while healthcare administrators can curate these assortments to maximize system-wide metrics. 

Our work studies a unique variant of assortment optimization with offline assortment and online response: assortments are created and offered simultaneously, and patients then respond sequentially in a random order to make provider selections. 
This variant is especially interesting because it captures the logistical complexities of the patient-provider matching problem. Dynamically varying assortments online are logistically taxing for healthcare administrators, yet it is difficult to control the patient response order, necessitating a combination of offline and online modeling. 
Moreover, this setting generalizes beyond patient-provider matching; similar situations can be found in other two-sided matching markets such as food delivery, where companies may propose a set of order options (menus) to drivers who respond in an online fashion, and where there is a tradeoff between driver autonomy and system utility (i.e., cost).

Our contributions answer this key question: \textit{How should healthcare administrators design provider menus for patients to optimize system-wide match rates and match quality?}
Our results show that limiting assortment sizes and offering carefully tailored menus enables a tradeoff between match rate and match quality, underscoring the balance between patient autonomy and centralized planning.

\subsection{Contributions}
Overall, we make three contributions to the patient-provider matching and assortment optimization literature: we i) model patient-provider matching using a hybrid offline-online variant of assortment optimization, ii) characterize the performance of assortment policies theoretically and empirically, and detail how the best-performing policy depends on problem characteristics such as patient match probability and patient/provider ratio, and iii) demonstrate that carefully tailoring assortment menus allows us to tradeoff between match rate and match quality.  

\paragraph{Patient-Provider Matching as Assortment Optimization}
We develop a model of patient-provider matching where administrators offer assortments upfront, and then patients respond sequentially in a random order. 
We do so to capture two critical elements: 1) patients have autonomy in selecting providers, and 2) the system has some control over the matching process in selecting assortments for each patient. 
Additionally, such a model can be broadly applicable to two-sided marketplaces where one side (users) selects online from an assortment, yet dynamic variation of the assortment is infeasible. 
We model patient decisions through a choice model that captures how the set of offered providers impacts the resultant matches. 
The selection of the choice model is key to the fidelity of our model, so we analyze various choice models, including the uniform choice model and the multinomial logit choice model.
Finally, we capture patient heterogeneity through a match quality matrix, which captures differences in match quality between pairs of patients and providers. 


\paragraph{Characterizing Assortment Policies}
We study various policies for assortment optimization and characterize their performance theoretically and empirically. 
We quantify policy performance through two metrics: match rate and match quality. 
We first propose a baseline greedy policy that offers all providers to all patients, then demonstrate that this policy can perform arbitrarily poorly for match quality. 
We then build upon this by developing new policies that adaptively vary assortment sizes.
We demonstrate that tailored assortment offerings can improve match quality compared to baseline greedy solutions. 
We empirically characterize the best policy as a function of model parameters such as the ratio of patients to providers and the underlying model assumptions. 
\footnote{We include all code at \url{https://github.com/naveenr414/patient-provider/}} 

\paragraph{Real-World Matching System Insights}
We apply our policies in a semi-synthetic simulation based on a real-world healthcare system in Connecticut. 
We demonstrate that the best-performing policy here outperforms the greedy baseline by 13\% for match quality. 
We extend this to consider metrics beyond match quality and match rate, such as fairness and regret. 
These findings yield actionable recommendations and highlight the tensions between various metrics when designing real-world patient-provider matching systems. 
Our key takeaway is that the best set of assortments offered depends on problem specifics, including the ratio of patients to providers and the willingness of patients to match. 
As patients are more willing to match, assortments should be smaller and more tailored to maximize match quality. 

\begin{figure*}
    \centering 
    \includegraphics[width=\textwidth]{figures/patient_provider.pdf}
    \caption{ Overview of patient-provider matching model. A) We match patients with providers B) by offering a set of assortments to patients, where each assortment consists of a set of providers. C) Patients then respond sequentially in a random order and either select a provider from their assortments or abstain. D) This results in matches between patients and providers.}
    \label{fig:pull}
\end{figure*}

\subsection{Related Work}
    \textbf{Matching Patients and Providers} 
Care continuity is critical for patient health, as it allows for better patient communication, lowers operating costs, and improves provider teamwork~\citep{lower_turnover}. 
Despite this, provider turnover rates are high, with the average healthcare system reporting a 7\% turnover rate year-to-year~\citep{lower_turnover}. 
Moreover, such issues are exacerbated recently due to high rates of provider burnout~\citep{provider_burnout} and the Covid-19 pandemic~\citep{provider_burnout_covid19}. 
When a provider leaves, it is difficult for patients to find a new one; only 54\% of patients find their new provider within one year, and 6\% fail to find a new provider even after three years~\citep{finding_new_provider}. 
High levels of provider turnover can worsen primary care because patients lack access to routine checkups which can reduce accountability, leading to consequences such as lower medication adherence~\citep{pcp_turnover,importance_primary_care,race_concordance_adherence,concordance_race_cardio}. 
Provider turnover itself impacts patient satisfaction and health outcomes, but the quality of the new match is also consequential in both dimensions. Nevertheless, it is difficult to quantify what constitutes a ``high-quality'' match. Patients care mainly about geographic proximity~\citep{patient_decisions}, with concordance along race~\citep{race_concordance}, gender~\citep{gender_concordance}, and language~\citep{language_concordance} playing a secondary role. 
Communication style alignment can improve match quality~\citep{personal_factor_concordance}. These factors build trust which translates to better care, such as leading patients to more truthfully report symptoms~\citep{race_concordance_lung}. 

\textbf{Algorithms for Patient-Provider Matching}
Existing work in patient-provider matching typically adopts either a one-shot matching framework or a genetic programming-based approach. 
Within the one-shot framework, linear programs are used to manage provider panel sizes~\citep{panel_size_lp}; deferred acceptance frameworks enable two-stage matching~\citep{two_stage_matching}; and scheduling algorithms model optimal rates of on-demand scheduling~\citep{on_demand_percent}. 
An alternative approach to patient-provider matching uses genetic programming, which learns matches over time through a fitness function. 
Within this framework, one line of work uses genetic programming to find fair matches between patients and providers~\citep{two_stage_matching}, while another approach uses genetic programming to balance workloads between providers~\citep{genetic_non_dominated}. 
In contrast to these approaches, we incorporate patient autonomy into the process, which motivates our assortment-based approach. 

\textbf{Assortment Optimization} 
Our work analyzes patient-provider matching through assortment optimization. In this setting, retailers construct assortments for customers, and customers make decisions based on offered options~\citep{assortment_optimization}. 
Retailers use models to capture customer choice behavior, such as the multinomial logit~\citep{assortment_mnl} and nested logit~\citep{assortment_nested_logit}, and generate product menus accordingly. 
Assortments are ubiquitous in e-commerce settings~\citep{e_commerce_assortment} but also employed in other domains such as school choice~\citep{assortment_school} and dating markets~\citep{assortment_dating}. Both offline and online variants have been studied. 
In offline assortment optimization, assortments are offered one-shot and customers make decisions simultaneously and independently~\citep{assortment_optimization}, while in online assortment optimization, customers arrive sequentially, and options are dynamically offered for each~\citep{assortment_online}. 
Other variants of assortment optimization include settings where system parameters are unknown~\citep{robust_assortment}, are dynamically learned~\citep{dynamic_assortment}, or have constraints~\citep{assortment_constraint}. 
Our setting resides between offline and online, as we offer assortments offline, but patients respond online. 
\section{Problem Formulation}
\label{sec:problem}

\subsection{Formal Model of Patient-Provider Matching}
\label{sec:model}
We introduce the problem of matching $N$ patients to $M$ providers. 
We present each patient $i$ with an assortment $\mathbf{X}_{i} \in \{0,1\}^{M}$ upfront.
Here, $X_{i,j}$ denotes whether patient $i$ has provider $j$ in their assortment. 
Patients then select providers from their assortment sequentially by responding in an order $\sigma$, and can abstain from provider selection. 
Here, $\sigma_{1} \in [N]$ represents the first patient in the order. 
We assume patient response times are i.i.d. so $\sigma$ is a uniformly random permutation of $[N]$; prior work demonstrates this can model real-world wait times~\citep{queueing_theory} and we relax this assumption in Section~\ref{sec:assumption_relaxation}. 

Patients select providers based on match quality $\theta_{i,j} \in [0,1]$, where $\theta_{i,j}$ is analogous to a reward for patient $i$ selecting provider $j$. 
Match quality can encompass various factors, including demographic concordance~\citep{race_concordance,gender_concordance,language_concordance}, physical proximity, and patient needs.
We assume administrators know $\theta_{i,j}$ because $\theta_{i,j}$ can be estimated through patient surveys; we further discuss this assumption in Section~\ref{sec:assumption_details}. 

Each patient selects a provider from their assortment according to a choice model $f_{i}: \{0,1\}^{M} \rightarrow \{0,1\}^{M}$. 
The choice model takes as input a 0-1 vector $\mathbf{z}$ denoting available and offered providers, and it outputs a random 0-1 vector representing the provider selected by patient $i$. 
Formally, let $\mathbf{y}^{(t)}$ represent the set of unselected providers at time $t$. 
Then, patient $\sigma_{t}$ selects providers according to $f_{\sigma_{t}}(\mathbf{X}_{\sigma_{t}} \odot \mathbf{y}^{(t)})$, where $\odot$ is the element-wise product ($\mathbf{a} \odot \mathbf{b} = \mathbf{c} \rightarrow a_{k} b_{k} = c_{k} \forall k$). 
$\mathbf{X}_{\sigma_{t}} \odot \mathbf{y}^{(t)}$ represents the set of providers that are on patient $\sigma_{t}$'s assortment and unselected by previous patients. 
If $f_{\sigma_{t}}(\mathbf{X}_{\sigma_{t}} \odot \mathbf{y}^{(t)}) = \mathbf{0}$, then no provider is selected, while if $f_{\sigma_{t}}(\mathbf{X}_{\sigma_{t}} \odot \mathbf{y}^{(t)})_{j} = 1$, then provider $j$ is selected. 
We note that a patient can match with at most one provider: $\lVert f_{\sigma_{t}}(\mathbf{X}_{\sigma_{t}} \odot \mathbf{y}^{(t)})\rVert_{1} \leq 1$. 
We assume that each provider can match with at most one patient as well. 
We assume this to simplify the presentation of our paper, though we relax this assumption in Section~\ref{sec:assumption_relaxation} and show in Appendix~\ref{sec:varied_capacity} that this is empirically equivalent to modifying the total number of providers.
We present examples of choice models below: 
\begin{enumerate}
    \item \textbf{Uniform} - Patients select their most preferred available provider with probability $p$ and otherwise abstain. 
    Let $\mathbf{e}_{k}$ denote the $k$-th standard basis vector, then: 
    
    \begin{equation}
    f_{i}(\mathbf{z}) =\left\{ \begin{array}{ll}
                \mathbf{e}_{\argmax_{j} (\theta_{i} \odot \mathbf{z})_{j}} & \text{with prb. p}\\
                \mathbf{0} &\text{otherwise}
            \end{array} \right.
    \end{equation}

    \item \textbf{Threshold} - Patients follow the uniform choice model but only select providers with match quality above some threshold $\alpha$. For example, if match quality corresponds to geographic proximity, $\alpha$ corresponds to maximum travel distance. 
Formally: 
    {\small \begin{equation}
    f_{i}(\mathbf{z}) =\left\{ \begin{array}{ll}
                \mathbf{e}_{\argmax_{j} (\theta_{i} \odot \mathbf{z})_{j}} & \text{with prb. p if } \max \theta_{i} \odot \mathbf{z} \geq \alpha \\
                \mathbf{0} &\text{otherwise}
            \end{array} \right.
    \end{equation}}
    \item \textbf{Multinomial Logit (MNL)} - Patients select providers according to $\theta$, and have an exit option valued at $\gamma$ which represents selecting no provider. Under the MNL choice model~\citep{choice_models_textbook}, $f$ is:
    \begin{equation}
            f_{i}(\mathbf{z}) =\left\{ \begin{array}{ll}
                \mathbf{e}_{j} & \text{with prb. } \frac{z_{j}\exp(\theta_{i,j})}{\exp(\gamma) + \sum_{j'} z_{j'} \exp(\theta_{i,j'})}\\\\
                \mathbf{0} &\text{with prb. } \frac{\exp(\gamma)} {\exp(\gamma) + \sum_{j'} z_{j'} \exp(\theta_{i,j'})}
            \end{array} \right.
    \end{equation}
\end{enumerate}
% Patients select providers according to a true match quality $\theta^{*}$, but we only have knowledge of $\theta_{i,j} = \theta^{*} + \epsilon$, where $\epsilon$ is Gumbel distributed noise~\citep{choice_models_textbook}. 
    % Patients either select a provider $j$ or can abstain from selecting any provider, which we represent through an exit option with $\gamma$: 

We aim to find a policy $\pi(\theta)$ that constructs an assortment $\mathbf{X}$ from a match quality matrix $\theta$. 
We evaluate policies through two metrics: match quality and match rate. 
\begin{enumerate}  
    \item \textbf{Match Rate (MR)} is the proportion of patients who are matched. Formally, it computes if patient $\sigma_{t}$ matches with a provider: $\lVert f_{\sigma_{t}}(\mathbf{X}_{\sigma_{t}} \odot \mathbf{y}^{(t)}) \rVert_{1} \geq 0$: 
    \begin{equation}
        \mathrm{MR}(\mathbf{X},\theta,f) = \mathbb{E}_{\sigma}[\frac{1}{N} \sum_{t=1}^{N} \lVert f_{\sigma_{t}}(\mathbf{X} \odot \mathbf{y}^{(t)}) \rVert_{1}]
    \end{equation}
    \item \textbf{Match Quality (MQ)} measures $\theta$ across selected patient-provider pairs. Formally, it aggregates $\theta_{\sigma_{t}}$ across provider selections $f_{\sigma_{t}}(\mathbf{X} \odot \mathbf{y}^{(t)})$ for all $N$ patients (= $N$ time periods):
    \begin{equation}
       \mathrm{MQ}(\mathbf{X},\theta,f) = \mathbb{E}_{\sigma}[\frac{1}{N} \sum_{t=1}^{N} f_{\sigma_{t}}(\mathbf{X} \odot \mathbf{y}^{(t)}) \cdot \theta_{\sigma_{t}}]
    \end{equation} 
\end{enumerate}
We select these two metrics because they naturally capture the scale and quality of matches, though, in Section~\ref{sec:real_world}, we evaluate additional metrics such as fairness and regret. 

\begin{example}
    We quantify the example in Figure~\ref{fig:pull} to demonstrate the metric calculations. 
    Suppose $N=M=3$, and let $\theta$ and $\mathbf{X}$ be 
    \begin{equation} 
        \theta =
        \begin{bmatrix}
        0.5 & 0.6 & 0.9 \\
        0.7 & 0.5 & 0.4 \\
        0.6 & 0.2 & 0.8
        \end{bmatrix}, \hspace{0.5cm}
        \mathbf{X} =
        \begin{bmatrix}
        0 & 0 & 1 \\
        1 & 1 & 0 \\
        1 & 0 & 1
        \end{bmatrix}
    \end{equation}
    Suppose $\sigma = [2,1,3]$. 
    Then one trial could be the following: patient $\sigma_{1}=2$ abstains from selecting a provider, patient $1$ selects provider $3$, and patient $3$ selects provider $1$: $f_{2}(\mathbf{y}^{(1)} \odot \mathbf{X}_{2}) = \mathbf{0}$, $f_{1}(\mathbf{y}^{(2)} \odot \mathbf{X}_{1}) = \mathbf{e}_{3}$, and $f_{3}(\mathbf{y}^{(3)} \odot \mathbf{X}_{3}) = \mathbf{e}_{1}$. 
    Here, $\mathbf{y}^{(1)} = [1,1,1]$, $\mathbf{y}^{(2)} = [1,1,1]$, and $\mathbf{y}^{(3)} = [1,1,0]$. The resulting match rate is $\frac{2}{3}$ and match quality is $\frac{1}{3}(\theta_{1,3} + \theta_{3,1}) = 0.5$.
\end{example}

\subsection{Relevance to Patient-Provider Matching}
\label{sec:assumption_details}
We provide additional details on the reasons behind the assumptions made in Section~\ref{sec:model} and their connections to realities in healthcare. 

\paragraph{Offline-Online Setting} 
We assume that patients are offered assortments offline and then respond in an online manner; this choice reflects logistical constraints and administrative burden.
While pure offline assignment limits autonomy, pure online assortments or dynamically modifying the assortments online results in two problems: a) potentially misaligned incentive for patients to wait to make decisions, and b) an additional logistical burden. 
If assortments are dynamically modified, then, for example, patients can see different providers depending on the day they log in to the patient portal (e.g., A-B-C on Monday, X-Y-Z on Wednesday), which leads to confusion and incentivizes waiting. 
Dynamic offering might also be infeasible if patient notifications are sent in an analog manner (e.g., snail mail), and offering assortments in an online fashion might require administrators to continually update and reorganize assortments.
In Section~\ref{sec:assumption_relaxation}, we show that relaxing the model to a fully online setting has little impact on our overall takeaways.

\paragraph{Assumptions} 
We make several key assumptions in the model definition: a) providers each have unit capacity, b) patients respond in a uniform random order, and c) the match quality $\theta$ is perfectly known. 
Our core analytical results rely on these assumptions, but we perform sensitivity analyses to study our policies' performance in more general settings. We detail how to relax the unit capacity assumption and provide results in Section~\ref{sec:assumption_relaxation} and demonstrate that this has little impact on our policies or findings. 
We assume that patients respond in random order to simplify the problem setup, and to better understand this assumption in practice, we experiment with non-uniform random orderings in Section~\ref{sec:assumption_relaxation}. 
Finally, while healthcare administrators do not currently explicitly estimate match quality $\theta$, this notion is anecdotally incorporated into rematching decisions, such as finding proximal providers based on a patient's location. To quantify match quality, administrators could collect preferences over providers for each patient and estimate rankings over providers via a learning-to-rank algorithm~\citep{learning_to_rank}; we leave this estimation problem to future study. 
Even if $\theta$ is not perfectly known, we demonstrate in Appendix~\ref{sec:misspecification} that estimates of $\theta$ suffice to allow for good performance. 

\paragraph{Model Fidelity} 
Our model aims to capture real-world concerns by balancing patient autonomy with system-wide utility through assortment offer decisions. 
The emphasis on patient autonomy arises from our conversations with healthcare partners, where patients might prefer a few high-quality choices to an overwhelmingly large number of options~\citep{choice_overload}, and from prior work that demonstrates people are willing to trade autonomy for efficiency, including in online markets~\citep{autonomy_efficiency} and domestic life~\citep{india_efficiency_autonomy}. 
Here, we intentionally leave the definition of ``quality'' ambiguous, as the definition of match quality should reflect patient preferences.  
For example, if patients primarily care about proximity, match quality should be based on the distance between patients and providers. 
We discuss this further in Section~\ref{sec:semi_synthetic_data}. 

\subsection{Analysis in the Single Provider Scenario}
\label{sec:greedy_provider}
To gain insight into the optimal policy for both match rate and match quality, we analyze the situation when $M=1$: 

\begin{restatable}{theorem2}{thmmone}
\label{thm:one_m}
    Let $f$ be the uniform choice model with probability $p$.
    Suppose $M=1$, and let $u_{1},u_{2},\ldots,u_{N}$ be a permutation of $\{1,\ldots,N\}$ such that $\theta_{u_{1},1} \geq \theta_{u_{2},1} \cdots \theta_{u_{N},1}$; that is patient $u_{1}$ has the highest match quality, followed by $u_{2}$, etc..
    Let $s$ be defined as follows:
    \begin{equation}
       s = \argmax_{s} (1-(1-p)^{s}) \frac{\sum_{i=1}^{s} \theta_{u_{i},1}}{s}
    \end{equation}
    Then the policy which maximizes match rate is $\pi(\theta) = \mathbf{X} = \mathbf{1}$ while the policy which maximizes match quality is $\pi(\theta) =  \mathbf{X}, \mathbf{X}_{u_{1},1} = \mathbf{X}_{u_{2},1} \cdots \mathbf{X}_{u_{s},1} = 1$, where $\mathbf{X}_{i,1}$ is 0 otherwise. 
\end{restatable}

We include all full proofs in Appendix~\ref{sec:proofs}, and sketch the brief intuition here. 
Here $s$ represents the number of patients we offer the single provider to, and $u$ represents the ordering of patient match qualities for a given provider. 
When $M=1$, we can decompose the match quality into the probability any patient matches and the average match quality.
When the single provider is offered to $s$ patients, the match rate is $ (1-(1-p)^{s})$, while the average match quality for the top $s$-patients is $\frac{\sum_{i=1}^{N} \theta_{u_{i},i}}{s}$. 
Therefore, to maximize match rate we let $s=N$, while for match quality, we find the $s$ which maximizes the product of the average match quality and the match rate. 
Our results demonstrate a tradeoff between match rate and match quality even in the single provider scenario, and we next detail algorithms to generalize this idea. 

\section{Constructing Assortment Policies}
\label{sec:policies}
We develop and analyze policies $\pi(\theta)$ to construct assortments that optimize for match rate and match quality. 
We first introduce a baseline greedy policy, which offers all providers to all patients. 
While the greedy policy maximizes the match rate, it can lead to poor match quality, motivating the need for more complex policies. 
As a result, we propose three new policies: pairwise, group-based, and gradient descent in Section~\ref{sec:introduce_policies} and analyze their match rate and match quality in Section~\ref{sec:match_rate} and Section~\ref{sec:match_quality} (summarized in Table~\ref{tab:policy_summary}).

We first introduce an illustrative example and use this throughout this Section. 
\begin{example}
    \label{ex:illustration}
    Consider a scenario with $p=0.75$ and $\theta = [0.7,0.7,0.1]$. 
    According to Theorem~\ref{thm:one_m}, the optimal assortment is $\mathbf{X} = [1,1,0]$, which achieves a match quality of $0.22$. 
\end{example}


\subsection{Assortment Policies}
\label{sec:introduce_policies}
\subsubsection{Greedy Policy}
\label{sec:greedy}
The \textbf{greedy policy} offers all providers to each patient and grants patients full autonomy. 
Formally, the greedy policy is $\pi^{R}(\theta)_{i,j} = 1$, for every patient $i$ and provider $j$. 
The greedy policy maximizes the match rate under the uniform choice model because $\mathbf{X}_{\sigma_{t}} = \mathbf{1}$, and so $\lVert f_{\sigma_{t}}(\mathbf{X}_{\sigma_{t}} \odot \mathbf{y}^{(t)})\rVert_{1} = \lVert f_{\sigma_{t}}(\mathbf{y}^{(t)}) \rVert_{1}$.

For match quality, the greedy policy considers assortments myopically and ignores the structure of the match quality matrix, which can lead to poor performance. 
For example, if $M=1$ and $\theta_{1,1} >> \theta_{2,1}$, the greedy policy gives both patients equal chances of matching, whereas the optimal policy offers $j=1$ only to patient $i=1$. 

\subsubsection{Pairwise Policy}
\label{sec:pairwise}
To motivate the need for policies beyond greedy, we demonstrate that the greedy policy is sub-optimal in Example~\ref{ex:illustration}.
Recall that $\theta = [0.7,0.7,0.1]$, so it is better for patients 1 or 2 to match rather than patient 3. 
The greedy policy has $\pi^{R}(\theta)  = [1,1,1]$, which leads to an average match quality of $\frac{(1-(1-p)^{N}) \frac{1}{N} \sum_{i=1}^{N} \theta_{i,1}}{N} = 0.16$. 
Greedy achieves a low match quality because it offers a provider to patient $3$, which drags down the average match quality.
More broadly, the greedy policy can offer providers to too many patients, which can lead to sub-optimal match qualities because assortments are not tailored. 

To better tailor assortments to patients, we introduce the \textbf{pairwise policy} $\pi^{P}(\theta)$, which contrasts with the greedy policy by offering each patient at most one provider. 
The pairwise policy, $\pi^{P}(\theta)$, pairs patients and providers by solving the weighted bipartite matching problem: 
\begin{equation}
\max\limits_{\substack{
\pi^{P}(\theta), \\
\sum_{i=1}^{N} \pi^{P}(\theta)_{i,j} \leq 1, \\
\sum_{j=1}^{M} \pi^{P}(\theta)_{i,j} \leq 1
}} 
\sum_{i=1}^{N} \sum_{j=1}^{M} \pi^{P}(\theta)_{i,j} \theta_{i,j}
\end{equation}
This strategy better incorporates match quality structure at the expense of fewer patient choices. 
We note that we can optimize $\pi^{P}(\theta)$ by solving a linear program. 

\subsubsection{Group-Based Policy}
\label{sec:grouping}
While the pairwise policy can potentially better tailor assortments to patients, it can suffer from assortments that are too constrained. 
For example, returning to Example~\ref{ex:illustration}, the pairwise policy offers $\mathbf{X} = [1,0,0]$, which achieves a match quality of $0.18$. 
While the pairwise policy improves upon the greedy policy for match quality, it is still not the optimal assortment ($\mathbf{X} = [1,1,0]$) because it only offers each provider at most one patient. 

To fix this problem, we design the \textbf{group-based} policy $\pi^{G}$, which expands assortments from the pairwise policy to improve match quality. 
The group-based policy heuristically groups together matches from the pairwise policy to form patient ``groups'' that are all offered the same providers. 
To define the group-based policy, first let $v(\theta)_{i} = j$ if $\pi^{P}(\theta)_{i,j} = 1$; that is, $v(\theta)_{i}$ denotes the matches from the pairwise policy, where $v(\theta)_{i} = -1$ if $\pi^{P}(\theta)_{i,j} = \ 0 \ \forall j$. 
The group-based policy proceeds in two steps: 1) weight computation and 2) group formation. 
1) We first compute edge weights $\alpha_{i,i'}$ for each pair of patients $i,i'$. 
Here, $\alpha_{i,i'}$ is an estimation of the benefit in match quality when $i$ and $i'$ are put into a group with an assortment that is the union of $\mathbf{X}_{i}$ and $\mathbf{X}_{i'}$. 
Formally, let $\mathbf{X}^{\prime} = \pi^{P}(\theta), {X}^{\prime}_{i,v_{i'}}= X^{\prime}_{i',v(\theta)_{i}}=1$, then  $\alpha_{i,i'} = \mathrm{MQ}(\mathbf{X}',\theta,f) - \mathrm{MQ} (\pi^{P}(\theta),\theta,f)$. 
2) We then repeatedly construct groups $\mathbf{q}$ that maximize the sum of edge weights $\alpha_{i,i'}$ within the group. 
That is we optimize: 
\begin{equation}
    \max\limits_{\mathbf{q}} \sum_{i} \sum_{i'>i} q_{i} q_{i'} \alpha_{i,i'}
\end{equation}
We note that we can find such a solution quickly by letting $z_{i,j} = q_{i} q_{j}$, noting that $q_{i} \in \{0,1\}$, and solving at most $N$ linear programs (though in practice, it requires far fewer computations). 
We repeatedly solve for $\mathbf{q}$ and construct assortments from this. 
By constructing assortments in this group-based format, we guarantee that the match rate of our group-based policy matches the pairwise policy; we provide details in Section~\ref{sec:match_rate}. 
We formalize this procedure in Algorithm~\ref{alg:grouping}. 

\begin{algorithm}[h]
   \caption{Group-based policy ($\pi^{G}$)}
\begin{algorithmic}[1]
    \STATE {\bfseries Input:} Match quality matrix $\theta$ 
    \STATE {\bfseries Output:} Assortment, $\mathbf{X}$
    \STATE Initialize the assortment $\mathbf{X} = \pi^{P}(\theta)$
    \FORALL{$(i,i') \subseteq [N]$}
        \STATE   $\mathbf{X}^{\prime} = \pi^{P}(\theta), {X}^{\prime}_{i,v_{i'}}= X^{\prime}_{i',v(\theta)_{i}}=1$ 
        \STATE Let $\alpha_{i,i'} = \mathrm{MQ}(\mathbf{X}',\theta,f) - \mathrm{MQ} (\pi^{P}(\theta),\theta,f)$ 
    \ENDFOR

    \STATE Let $s=\infty$ and $r=\{1,\ldots,N\}$
    \WHILE{$s > 0$}
        \STATE Let $s = \max\limits_{\mathbf{q}} \sum_{i \in r} \sum_{\substack{i' \in r \\ i' > i}} q_{i} q_{i'} \alpha_{i,i'}$, $\mathbf{q}^* = \arg\max\limits_{\mathbf{q}} \sum_{i \in r} \sum_{\substack{i' \in r \\ i' > i}} q_i q_{i'} \alpha_{i,i'},$
        \STATE Let $X_{i,v_{i'}} = 1$ for all $i,i'$ with $q^*_{i} = q^*_{i'} = 1$ 
        \STATE Remove all $i$ with $q^*_{i} = 1$ from $r$
    \ENDWHILE
\end{algorithmic}
\label{alg:grouping}
\end{algorithm}

\subsubsection{Gradient Descent Policy}
\label{sec:lower_bound}
We finally propose the \textbf{gradient descent} policy $\pi^{D}$, which generalizes beyond pairwise interactions by optimizing for a surrogate heuristic for match quality under the uniform choice model. 

The heuristic consists of two functions: $h(\mathbf{X}) \in [0,1]^{N \times M}$, where $h(\mathbf{X})_{i,j}$ denotes the probability patient $i$ has provider $j$ available in their assortment, and $g(h(\mathbf{X})) \in [0,1]^{N \times M}$, where $g(h(\mathbf{X}))_{i,j}$ denotes the probability that provider $j$ is the most-preferred available provider for patient $i$. 
Here, $g$ and $h$ are known and defined below, so we aim to find a $\mathbf{X}$ to optimize $p \cdot \langle g(h(\mathbf{X})),\theta \rangle$. 
Under the uniform choice model, the probability that patient $i$ matches with provider $j$ is then $p \cdot g(h(\mathbf{X}))$, and so the expected match quality is $p \cdot \langle g(h(\mathbf{X})), \theta \rangle$. 

To compute $h(\mathbf{X})$ and $g(h(\mathbf{X}))$,
first recall that $\sigma$ is a random ordering of patients, and $\sigma_{t} = i$ denotes that patient $i$ is the t-th patient to select their provider. 
To find $h(\mathbf{X})_{i,j}$, we multiply $X_{i,j}$, which indicates if $j$ is offered to $i$, by the probability that $j$ is still available when patient $i$ arrives. 
We approximate the latter by considering the availability probability for each fixed arrival point (i.e., if $\sigma_t = i$) and average over all possible $N$ arrival orders. 
Clearly, if $\sigma_{1} = i$, then $h(\mathbf{X})_{i,j} = X_{i,j}$, as no previous patient could have selected provider $j$. 
If $\sigma_t = i$, then the probability that none of the previous $t-1$ patients select provider $j$ is at least $(1-p)^{t-1}$; at most $t-1$ prior patients could have been offered $j$ and had $j$ as their most preferred provider. 
Because $t$ is uniformly distributed from $1$ to $N$: 
\begin{equation}
    h(\mathbf{X})_{i,j} = X_{i,j}\sum_{t=1}^{N} \frac{(1-p)^{t-1}}{N} 
\end{equation}
We can further tighten this by considering the number of patients who are offered provider $j$, namely $n = \lVert \mathbf{X}_{*,j} \rVert_{1}$ out of $N$ patients. Since only $n$ patients are offered provider $j$, we rescale the selection probabilities accordingly: $p' = \left(\frac{n-1}{N-1}\right)p$. 
We can then update our calculation of $h(\mathbf{X})_{i,j}$ as:
{\small \begin{equation}
h(\mathbf{X})_{i,j} = X_{i,j}\sum_{t=1}^{N} \frac{(1-p'_{j})^{t-1}}{N}\  \mathrm{ where } \ p'_{j} = \left( \frac{\lVert \mathbf{X}_{*,j} \rVert_{1}-1}{N-1} \right) p
\end{equation}}

We then compute $g(h(\mathbf{X}))$ by first computing a preference ordering over providers for patient $i$. 
Let $u_{i,1},u_{i,2},\ldots,u_{i,M}$ be the most to least preferred providers for patient $i$; that is, $\theta_{i,u_{i,1}} \geq \theta_{i,u_{i,2}} \geq \cdots \geq \theta_{i,u_{i,M}}$. 
Then for patient $i$ to pick their kth best option, it requires that 1) $u_{i,k}$ is available, and 2) none of $u_{i,1},u_{i,2},\ldots,u_{i,k-1}$ are available. 
Formally, we can estimate the probability that the kth most preferred provider is the best available as: 
\begin{equation}
    g(h(\mathbf{X}))_{i,u_{i,k}} = h(\mathbf{X})_{i,u_{i,k}} \prod_{k'=1}^{k-1} (1-h(\mathbf{X})_{i,u_{i,k'}})
\end{equation}
Using $h$ and $g$, we can compute a differentiable heuristic and in Section~\ref{sec:gradient_descent_match_quality}, we demonstrate that this heuristic lower bounds the total match quality.  
In practice, we use gradient descent and include a log regularization term to ensure that $\mathbf{X}$ is binary. 
While the heuristic is not always convex, in Section~\ref{sec:empirical}, we demonstrate that the policy is still able to find well-performing solutions, implying that a well-performing minimum is found despite guarantees on convexity. 
Moreover, such solutions can be found quickly in practice because the gradient descent algorithm quickly converges to local minima (see Section~\ref{sec:empirical}). 

\begin{algorithm}[h]
   \caption{Gradient Descent Objective}
\begin{algorithmic}
    \STATE {\bfseries Input:} Match quality matrix $\theta$ and Assortment $\mathbf{X}$
    \STATE {\bfseries Output:} Estimated total match quality
    \STATE Let $p'_{j} = \frac{\lVert \mathbf{X}_{*,j} \rVert_{1}-1}{N-1}$
    \STATE Let $h(\mathbf{X})_{i,j} = \frac{X_{i,j}}{N} \sum_{t=1}^{N} (1-p'_{j})^{t-1}$ 
    \STATE Let $u_{i,1},u_{i,2},\ldots,u_{i,M}$ be a permutation of $[M]$, so $\theta_{i,u_{i,1}} \geq \theta_{i,u_{i,1}} \cdots \theta_{i,u_{i,M}}$
    \STATE Let $g(h(\mathbf{X}))_{i,u_{i,k}} = h(\mathbf{X})_{i,u_{i,k}} \prod_{k'=1}^{k-1} (1-h(\mathbf{X})_{i,u_{i,k'}})$
    \RETURN $p \cdot \langle g(h(\mathbf{X})), \theta \rangle$
\end{algorithmic}
\label{alg:provider}
\end{algorithm}

\subsection{Match Rate Analysis}
\label{sec:match_rate}

\begin{table}[tb]
\centering
\caption{We prove lower bounds for the match rate (MR) and match quality (MQ) for all four policies. 
We bold all tight bounds. 
The main takeaway is that the greedy policy maximizes match rate, but is poor for match quality, incentivizing the need for better policies.}
\begin{tabular}{@{}lccl@{}}
\toprule
Policy       & \multicolumn{1}{l}{Match Rate (MR)}    & \multicolumn{1}{l}{Match Quality (MQ)} & Notes                                             \\ \midrule
Greedy           & $\geq \min(p,\frac{M}{N})$     & $< \boldsymbol{\epsilon} \ \forall\  \epsilon$                        & -                                \\
Pairwise         & $ = \mathbf{p \frac{\min(M,N)}{N}}$           & $\geq  \mathbf{p \mathrm{MQ}(\pi^{*}(\theta),\theta,f)}$                               &            -                                       \\
Group-Based      & $= \mathbf{p \frac{\min(M,N)}{N}}$           & -                                 & \makecell[l]{Heuristically improves MQ \\ vs. pairwise} \\
Gradient Descent & $< p \frac{\min(M,N)}{N}$ & -                                 & \makecell[l]{Provably optimizes lower \\ bound of MQ} \\ \bottomrule
\end{tabular}
\label{tab:policy_summary}
\end{table}

We formally characterize the match rate for the four policies, and find that the greedy policy theoretically performs best, while gradient descent performs worst. 
We focus on the uniform choice model due to its simplicity and ability to model the decision-making of real-world patients~\citep{choosing_doctor}. 
We include proof sketches here, while full proofs are deferred to Appendix~\ref{sec:proofs}. 

\subsubsection{Greedy Policy}
We first characterize the performance of the greedy policy: 

\begin{restatable}{theorem2}{thmgreedymatch}
\label{thm:greedy_match}
    Let $f$ be the uniform choice model with match probability $p$. 
    Then 
    \begin{equation}
        \mathrm{MR}(\pi^{R}(\theta),\theta,f)  \geq \min(p,\frac{M}{N})
    \end{equation}
    Moreover, there exists no policy $\pi'$ such that $\mathrm{MR}(\pi'(\theta),\theta,f) > \mathrm{MR}(\pi^{R}(\theta),\theta,f)$
\end{restatable}

We demonstrate this by considering two scenarios; one in which all providers are matched and the other in which each patient gets a chance to match. 
Under the uniform choice model, the greedy policy is optimal for the match rate because it offers the largest assortments.   

\subsubsection{Pairwise Policy}
We similarly characterize the match rate of the pairwise policy: 

\begin{restatable}{theorem2}{thmpairwisematch}
\label{thm:pairwise_match}
     Let $f$ be the uniform choice model with match probability $p$. 
    Then 
    \begin{equation}
        \mathrm{MR}(\pi^{P}(\theta),\theta,f)  = p \frac{\min(M,N)}{N}
    \end{equation}
\end{restatable}
Here, we prove this by noting that assortments are disjoint, and each patient with a non-empty assortment matches with probability $p$. 
Noticing that there are $\min(M,N)$ non-empty assortments yields the desired match rate. 
We note that this bound is tight due to the disjoint nature of patient assortments. 
When $M \geq N$, the pairwise policy achieves the optimal match rate, while when $M \leq N$, the pairwise policy performs worse than optimal because the same provider should be offered to multiple patients. 

\subsubsection{Group-Based Policy}
\label{sec:group_match_rate}
Because we construct the group-based policy so assortments are shared across patients in a group, we can show that the group-based policy maintains the same match rate as the pairwise policy:

\begin{restatable}{theorem2}{thmgroupingmatch}
    \label{thm:grouping_match}
    Let $f$ be the uniform choice model. Then 
    \begin{equation}
        \mathrm{MR}(\pi^{G}(\theta),\theta,f) = \mathrm{MR}(\pi^{P}(\theta),\theta,f)
    \end{equation}    
\end{restatable}

We prove this by showing that the number of patients with an empty assortment is the same in both the group-based and pairwise policies, leading to the same number of matches in both policies. 
Providing structure, as we do for the group-based policy, is necessary to ensure match rates; once this structure is abandoned, it becomes difficult to ensure high match rates. 

\subsubsection{Gradient Descent Policy}
We characterize the match rate for the gradient descent policy in two steps. 
First, we demonstrate that there exist scenarios where certain maximums for $g(h(\mathbf{X}))$ lead to a poor match rate.
Next, we show that even when the maximum is unique, it can still lead to a worse match rate compared with the pairwise policy. 
Taken together, these results demonstrate that optimizing the match quality heuristic can lead to tensions with the match rate. 
\begin{example}
    Consider $\theta_{i,j} = 0$ except at $\theta_{1,1} = 1$. 
    Then for any assortment $\mathbf{X}$, only $g(h(\mathbf{X}))_{1,1}$ contributes to the match quality. 
    Therefore, to maximize $g(h(\mathbf{X}))_{1,1}$, we set $\mathbf{X}_{1,1} = 1$ and $\mathbf{X}_{i,1} = 0 \ \forall \ i$. 
    We note that the remaining $\mathbf{X}$ can be set arbitrarily.
    Therefore, letting $\mathbf{X}_{1,1} = 1$ and $\mathbf{X}_{i,j} = 0$ otherwise yields match rates as low as $\frac{1}{N}$. 
\end{example}
We next show that even when solutions are unique, the gradient descent policy can still lead to sub-optimal match rates: 

\begin{restatable}{theorem2}{thmgradientmatch}
    \label{thm:gradient_match}
    Let $f$ be the uniform choice model with match probability $p$. Then for any $p<1$, there exists $\theta$ such that the unique maximizer of $\langle g(h(\mathbf{X})),\theta \rangle$, $\mathbf{X}^{*}$, has: 

    \begin{equation}
        \mathrm{MR}(\mathbf{X}^{*},\theta,f) < \mathrm{MR}(\pi^{P}(\theta),\theta,f)
    \end{equation}
\end{restatable}

We prove this by constructing a simple $N=M=2$ example where even the unique maximizer can achieve a sub-optimal match rate. 
Here, the unique maximizer also corresponds to the true maximizer for match quality, and so we demonstrate a match rate-match quality tradeoff. 
% Finally, to give intuition for why the match rate might not be so bad in practice, we demonstrate bounds on the \textbf{ground-truth} function $g^{*}(h^{*}(\mathbf{X}))$. 
% Here, $g^{*}(h^{*}(\mathbf{X}))_{i,j}$ is the true probability that provider $j$ is the most preferred available provider for patient $i$, where the probability is taken over orderings $\sigma$ and provider availabilities $\mathbf{y}$.  
% \begin{restatable}{theorem2}{thmgradientmatchlowerbound}
%     \label{thm:gradient_match_lower_bound}
%     \textbf{Theorem~\ref{thm:gradient_match_lower_bound}}: Let $f$ be the uniform choice model with match probability $p$. 
%     Let $g^{*}(h^{*}(\mathbf{X}))_{i,j} = \mathrm{Pr}[\argmax \ \theta_{i} \odot \mathbf{y}^{(\sigma^{-1}_{i})} = j]$. 
%     Then there always exists some maximizer $\mathbf{X}^{*}$ of $p \cdot \langle g^{*}(h^{*}(\mathbf{X})),\theta \rangle$, such that: 
%     \begin{equation}
%     \mathrm{MR}(\mathbf{X}^{*},\theta,f) \geq \frac{1}{2} \mathrm{MR}(\pi^{P}(\theta),\theta,f) 
%     \end{equation}
% \end{restatable}
% We demonstrate this by demonstrating two properties: there always exists a maximizer so all providers are offered to some patient, and if a patient receives a large assortment, then each provider in their assortment must also be offered to other patients. 
% These two properties ensure that many patients have a non-empty assortment and increase their match probability. 
% As a result, we conjecture that the heuristic function $\langle g(h(\mathbf{X})),\theta \rangle$ has similar guarantees for \textit{some} maximizer, though the challenge in practice is finding such a maximizer when solution multiplicity exists. 
Overall, we find a tension between policies that optimize for match rate, such as greedy, and those that optimize for match quality, such as gradient descent. 

\subsection{Match Quality Analysis}
\label{sec:match_quality}
We next complement our match rate analysis by analyzing the match quality for various policies.
We show that the greedy policy, which performs best for match rate, can perform arbitrarily poorly for match quality.
We then demonstrate that our three new policies can improve match quality, showing a tradeoff between match rate and match quality. 

\subsubsection{Greedy Policy}
We start by demonstrating that the greedy policy $\pi^{R}$ can perform poorly when compared with the optimal match rate policy $\pi^{*}$. 
To demonstrate this, we first define $\pi^{*}$: 
\begin{equation}
    \pi^{*}(\theta) = \argmax_{\mathbf{X}} \mathrm{MQ}(\mathbf{X},\theta,f)
\end{equation}
We then demonstrate that $\pi^{R}$ is an $\epsilon$ approximation:

\begin{restatable}{theorem2}{thmgreedy}
    \label{thm:greedy}
    Let $f$ be the uniform choice model with match probability $p$. 
    For any $p$ and $\epsilon$, there exists a $\theta$ such that 
    \begin{equation}
        \mathrm{MQ}(\pi^{R}(\theta),\theta,f)  \leq \epsilon \mathrm{MQ}(\pi^{*}(\theta),\theta,f)
    \end{equation}
\end{restatable}

We prove this by generalizing an example where $\theta_{1,1} >> \theta_{2,1}$ to arbitrary $N$ and $M$. 
Here, the optimal policy is to offer each provider to only one patient. 
However, the greedy policy offers all providers to all patients, which leads to sub-optimal matches for each provider $j$. 

\subsubsection{Pairwise Policy}
We next show that the pairwise improves the performance guarantee from $\epsilon$ to $p$ for match quality: 

\begin{restatable}{theorem2}{thmlp}
\label{thm:lp}
    Let $f$ be the uniform choice model with match probability $p$. 
    Then 
    \begin{equation}
        \mathrm{MQ}(\pi^{P}(\theta),\theta,f) \geq p \mathrm{MQ}(\pi^{*}(\theta),\theta,f) 
    \end{equation}
\end{restatable}

To prove this, we first upper bound the optimal policy, $\pi^{*}(\theta)$, with the value of the underlying bipartite matching problem. 
We then show that the pairwise policy corresponds to the bipartite matching problem, where edges correspond to pairs of patients and providers, and each edge exists with probability $p$. 
Summing over edges gives that the pairwise policy is a $p$-approximation to the bipartite matching problem and is, therefore, a $p$-approximation to $\pi^{*}$. 
We note that when $M=1$ and $N$ is large, with $\theta_{i,1}=1 \ \forall \ i$, this bound becomes tight, as the optimal match quality is $1$, while pairwise achieves $p$. 

\subsubsection{Group-Based Policy}
We next analyze the match quality for the group-based policy. 
While proving a bound on the match quality for the group-based policy is combinatorially difficult, we give intuition for when a group-based policy is necessary by demonstrating the need to modulate between the greedy and pairwise policies:

\begin{restatable}{proposition2}{thmgrouping}
\label{thm:grouping}
Let $f$ be the uniform choice model with match probability $p$. 
If $\theta_{i,j} \sim U(0,1)$ and $M \leq N$ then
\begin{equation}
    \frac{\mathbb{E}_{\theta}[\mathrm{MQ}(\pi^{R}(\theta),\theta,f)]}{\mathbb{E}_{\theta}[\mathrm{MQ}(\pi^{P}(\theta),\theta,f)]}\geq \frac{1-(1-p)^{N/M}}{2p}
\end{equation}
\end{restatable}

When $p$ is small, the greedy policy can outperform pairwise because providers should be offered to multiple patients.
Meanwhile, larger $p$ incentivizes smaller menus, leading to better performance for the pairwise policy. 
Additionally, when $N>M$ and $p<\frac{1}{2}$, the greedy policy can outperform pairwise due to provider scarcity. 
The group-based policy is built on this intuition so that groups are dynamically sized as needed. 

\subsubsection{Gradient Descent Policy}
\label{sec:gradient_descent_match_quality}
To analyze the match quality of the gradient descent policy, we demonstrate that our heuristic $p \cdot \langle g(h(\mathbf{X})), \theta \rangle$ lower bounds the match quality: 

\begin{restatable}{theorem2}{thmlowerbound}
    \label{thm:lower_bound}
    The following holds for any $\mathbf{X}$ when $f$ is the uniform choice model with probability $p$:
    {\small \begin{equation}
        p \cdot \langle g(h(\mathbf{X})), \theta \rangle \leq \mathrm{MQ}(\mathbf{X},\theta,f)
    \end{equation}}
\end{restatable}

We demonstrate this by showing that $\sum_{j} g(h(\mathbf{X}))_{i,j}$ is an underestimate for any fixed $i$. 
Our lower bound becomes tighter as $N$ becomes larger than $M$ because there is less overlap between different providers' assortments due to scarcity. 

\subsection{Analysis under Other Choice Models}
\label{sec:beyond_uniform_choice}
Our analysis thus far has been primarily under the uniform choice model, so we discuss policy performance under the MNL choice model with an exit option $\gamma$, a common choice model seen throughout the assortment optimization literature~\citep{assortment_mnl}. 
We characterize the match rate and match quality for the pairwise policy under this model: 

\begin{restatable}{lemma2}{thmmnlmatch}
\label{thm:mnlmatch}
Let $f$ be the MNL choice model parametrized by $\gamma$. 
Then the match rate and match quality of the pairwise policy are: 
\begin{align}
        \mathrm{MR}(\pi^{P},\theta,f)  \\ = \frac{1}{N\sum_{i} \mathbbm{1}[v(\theta)_{i} \geq 0]} \sum_{i; v(\theta)_{i} \geq 0}^{} \frac{\exp(\theta_{i,v(\theta)_{i}})}{\exp(\theta_{i,v(\theta)_{i}}) + \exp(\gamma)}
\end{align}
and 
\begin{align}
        \mathrm{MQ}(\pi^{P}(\theta),\theta,f) \\  \geq \left(\frac{1}{N} \sum_{i=1}^{N} \frac{\exp(\theta_{i,v(\theta)_{i}})}{\exp(\theta_{i,v(\theta)_{i}}) + \exp(\gamma)}\right) \mathrm{MQ}(\pi^{*}(\theta),\theta,f)
\end{align}
\end{restatable}
To prove this, we apply similar techniques to those used with the uniform choice model by analyzing the individual match rate for each patient, and then summing across patients. We observe that the exit option parameter $\gamma$ plays a similar but inverted role to $p$ in the uniform choice model. 
As $\gamma$ increases, the probability a patient does not select any provider also increases, and so larger $\gamma$ incentivizes policies to offer larger assortments, similar to the incentives under smaller $p$ (see Proposition~\ref{thm:grouping}). 
We also note the additional complexity inherent to the MNL choice model. 
Guarantees under the MNL model require additional knowledge of $\theta$ and $v$, and analysis under the MNL choice model becomes difficult because match rates cannot be decoupled from individual match qualities. 
Nonetheless, the key takeaway is that a patient's willingness to match, here controlled by $\gamma$, can dictate policy performance. 
\section{Empirical Analysis}
\label{sec:empirical}
We complement our theoretical analysis through a study of our policies on a synthetic dataset. 
We characterize the best-performing policy based on model parameters, and show that the choice of assortment policy impacts match quality and match rate. 
Our biggest takeaway is that tailoring policies through gradient descent maximizes match quality, especially when patients outnumber providers. 

\subsection{Experimental Details}
We compare the four policies from Section~\ref{sec:policies} (greedy, pairwise, group-based, and gradient descent) and a random baseline with $\pi(\theta)_{i,j} \sim \text{Ber}(\frac{1}{2})$. 
For a fixed problem setup, we randomly sample $T$ permutations of patients and compute the match quality and match rate for the corresponding patient order ($\sigma$). We average runs across 15 seeds and $T=100$ trials and plot the standard error across these seeds. 
Because the scale of match quality and rate might vary between experiments, we report the normalized match quality and match rate (norm. MQ and norm. MR for short), which normalizes each quantity by that of the random policy.
All policies run in under five minutes, demonstrating the scalability of our policies across problem instances. 

\subsection{Impact of Model Parameters}
\label{sec:parameters}
We construct synthetic examples to gain insight into our policies across different scenarios. 
For each example, we specify values for $M$, $N$, $f$, and $\theta$. 

% \begin{example}
% \label{ex:small}
% \textbf{Few patients and providers} - We first analyze policies with a small number of patients $N$ and number of providers $M$. 
% Small $N$ and $M$ allow us to compare to the optimal policy $\pi^{*}$ because we can compute $\pi^{*}$ by enumerating all $2^{NM}$ assortments and selecting the one that maximizes match quality across samples of $\sigma$. 
% We run two experiments: the first fixes $N=2$ while varying $M$ from $2$ to $5$, and the second fixes $M=2$ while varying $N$ from $2$ to $5$. 
% For both, we uniformly distribute $\theta$, $\theta_{i,j} \sim U(0,1)$, and use the uniform choice model with $p=1/2$.
% We compute approximation ratios against the optimal $\pi^{*}$ in match quality. 

% \begin{figure*}
%     \centering 
%     \includegraphics[width=\textwidth]{figures/small_patient_provider.pdf}
%     \caption{We find the optimal policy and compute the approximation ratio for match quality compared to optimal. While all non-random policies achieve approximation ratios $\geq 0.9$ when $M \geq N$ (left), the pairwise policy performs poorly and achieves a low approximation ratio when $N \geq M$ (right).\hwcomment{extremely nit - same font size as Fig 3? looks big here}} 
%     \label{fig:small}
% \end{figure*}

% In Figure~\ref{fig:small}, we demonstrate that when $M \geq N$ (left), all non-random policies perform well and achieve within 3.5\% of optimal, while for $N \geq M$ (right), the pairwise policy performs poorly. 
% When $M \geq N$, providers are abundant so each patient can match with their most preferred provider. 
% However, when $N \geq M$, the pairwise policy performs up to 28\% worse than optimal. Providers are scarce in this scenario, so policies should offer the same provider to multiple patients. 
% The pairwise policy only offers each provider one patient, which is suboptimal and leads to poor match quality.
% \end{example}

\begin{example}
\label{ex:patient_provider}
\textbf{Impact of $N$ and $M$} - 
We first analyze how the patient/provider ratio impacts policy performance. 
We fix the number of providers $M=25$, while varying the number of patients $N$ from $25$ up to $200$. This reflects real-world situations in which patients generally outnumber providers~\citep{healthcare_shortage}.
We compare policies under two match quality distributions: the first uniformly distributes $\theta \sim U(0,1)$, while the second normally distributes $\theta_{i,j} \sim \mathcal{N}(\mu_{j},0.1^{2})$, where $\mu_{j} \sim U(0,1)$. 
The former corresponds to heterogeneous preferences where patient preferences are non-correlated, while the latter corresponds to homogeneous preferences, where patient preferences are influenced by some signal of provider quality $\mu_{j}$. 
For both experiments, we fix $p=0.5$ for the uniform choice model. 

\begin{figure*}[ht!]
    \centering 
    \includegraphics[width=\textwidth]{figures/vary_providers_patients_matches.pdf}
    \caption{When there are more patients than providers, the gradient descent policy performs nearly as well as the optimal greedy policy for match rate (top), and performs the best for match quality (bottom), for both uniformly and normally distributed $\theta$.}
    \label{fig:patient_providers}
\end{figure*}

In Figure~\ref{fig:patient_providers}, we demonstrate that the gradient descent policy performs best when patients outnumber providers. 
When $N/M=8$, the gradient descent policy performs 5\% better than alternatives for uniform $\theta$ and 9\% for normal $\theta$. 
The gradient descent policy performs best for large $N/M$ because $g(h(\mathbf{X}))$ better approximates match quality when $N>M$.
Both greedy and group-based policies perform similarly, while the pairwise policy performs worst. 
When comparing across match rates, we find that the greedy policy performs best as expected, while the gradient descent policy is nearly as good for $N/M \geq 4$. 
\end{example}

\begin{example}
\label{ex:comparison}
\textbf{Impact of $p$ and $\theta$} - 
\begin{figure*}[t]
    \centering 
    \subfloat[\centering Uniform Match Quality]{\includegraphics[height=0.25 \textwidth]{figures/phase_diagram.pdf}}      
    \hspace{2em}
    \subfloat[\centering Uniform Match Rate]{\includegraphics[height=0.25 \textwidth]{figures/phase_diagram_match.pdf}} \\
    \subfloat[\centering Normal Match Quality]{\includegraphics[height=0.25 \textwidth]{figures/phase_diagram_normal.pdf}}
    \hspace{2em}
    \subfloat[\centering Normal Match Rate]{\includegraphics[height=0.25 \textwidth]{figures/phase_diagram_normal_matches.pdf}} 

    \caption{We characterize the best-performing policy for match rate and match quality when varying $p$ and $\theta$. Both when $\theta$ is uniformly (top) or normally (bottom) distributed, we find that the gradient descent policy performs best for large $N/M$ and large $p$. Greedy policy performs best in match rate for small $p$ because larger assortments should be offered then. }
    \label{fig:phase}
\end{figure*}
To understand which policy performs best across $p$, $\theta$, and $N/M$, we plot the policy which maximizes match quality when varying $p \in \{0.1,0.3,0.5,0.7,0.9\}$ and $N/M \in \{1,2,3,4,5,6,7,8\}$ with $M=25$, and distributing $\theta$ uniformly (Figure~\ref{fig:phase} top) and normally (Figure~\ref{fig:phase} bottom).
We find that the gradient descent policy performs best for large patient/provider ratios, both in match rate and match quality, which matches the results from Example~\ref{ex:patient_provider}. 
The group-based policy performs well for small $p$ because the problem more closely resembles bipartite matching (see Theorem~\ref{thm:lp}), while the greedy policy performs well for small $p$ because it becomes advantageous to offer larger assortments (see Proposition~\ref{thm:grouping}).
When comparing across match rates, we find that the gradient descent policy is within 10\% of the greedy policy, and for larger $N/M$ or $p$ values, the gradient descent policy has the same match rate as greedy. 

\begin{figure*}
    \centering 
    \includegraphics[width=\textwidth]{figures/heterogeneity.pdf}
    \caption{We let $\theta$ be distributed according to $\mathcal{N}(\mu_{j},s^{2})$ and analyze the impact of the standard deviation $s$ on policy performance. We find that at smaller values of $s$, policies perform more similarly, while policy performance is more differentiated with more heterogeneous preferences.}
    \label{fig:heterogeneity}
\end{figure*}
To further characterize the impact of $\theta$ upon policy performance, we let $\theta$ be normally distributed according to $\mathcal{N}(\mu_{j},s^{2})$, and vary $s \in [0,0.25]$ while $M=25$ and $N=100$. 
Here, smaller $s$ correspond to homogeneous preferences across patients, as $\theta$ values have a smaller spread between patients for a fixed provider. 
In Figure~\ref{fig:heterogeneity}, we find that as $s$ increases, the gradient descent policy performs better, with gradient descent outperforming other policies by 10\% for $s=0.25$. 
When $s=0$, all patients have similar preferences, and so the exact selection of assortment makes little difference.    
However, as $s$ is increased, the particular pairings of patients and providers become more important, leading to improved performance for policies such as gradient descent which tailors policies for patients. 
\end{example}

\begin{example}
\label{ex:choice_model}
    \textbf{Impact of Choice Model $f$} - 
We evaluate the impact of choice models on policy performance by investigating the MNL and threshold choice models. 
We fix $M=25$ and $N=100$, let $\theta$ be distributed uniformly, and vary $\gamma \in \{0.1,0.25,0.5,0.75\}$ for the MNL choice model and vary  $\alpha \in \{0.1,0.25,0.5,0.75\}$ with $p=0.5$ for the threshold choice model.

\begin{figure*}
    \centering 
    \includegraphics[width=\textwidth]{figures/other_choice_models.pdf}
    \caption{We compare policy performances when varying the exit option $\gamma$ for the MNL choice model, and the threshold $\alpha$ for the threshold choice model. 
    For the MNL choice model, the choice of $\gamma$ dictates the performance of the pairwise policy as $\gamma$ controls a patient's willingness to match. For the threshold choice model, larger $\alpha$ leads to a bigger gap between gradient descent and other policies because gradient descent is better able to find high-quality matches. }
    \label{fig:other_choice}
\end{figure*}

In Figure~\ref{fig:other_choice}, we show that $\gamma$ determines the performance of the pairwise policy for the MNL choice model, while gradient descent performs better for larger $\alpha$ in the threshold choice model. 
As noted in Section~\ref{sec:pairwise}, the choice of $\gamma$ is inversely related to the setting of $p$, with small $\gamma$ analogous to large $p$. 
Therefore, smaller $\gamma$ leads to better guarantees for the pairwise policy. 
Comparing the greedy and pairwise policies, we find that at small $\gamma$, it is better to offer more tailored or smaller assortments, while for larger $\gamma$, it becomes better to offer larger assortments. 
For the threshold choice model, $\alpha$ impacts the performance of the gradient descent policy.
At $\alpha=0.1$, gradient descent outperforms alternatives by 4\%, while at $\alpha=0.75$, gradient descent is 11\% than all alternatives. 
When $\alpha$ is large, the gap between gradient descent and other policies increases because gradient descent can potentially guarantee higher quality matches when compared to policies such as greedy. 
\end{example}

\subsection{Relaxing Model Assumptions}
\label{sec:assumption_relaxation}
We relax several assumptions from Section~\ref{sec:model}, and assess the impact on policy performance: 
\begin{enumerate}
    \item \textbf{Dynamic Assortments} - Throughout the paper, we assume that assortments are initially computed offline and offered online due to the logistical burden of varying assortments dynamically. To understand the loss in performance due to this, we construct a simple ``dynamic'' heuristic, that re-computes the Pairwise policy but accounts for providers and patients who have already matched or declined matching. In essence, this allows assortments to update after each decision made by a patient. We note that our dynamic assortment policy is just one selection and is not necessarily optimal, and we leave it to future work to further explore dynamic policies. 
    \item \textbf{Non-Uniform Random Response} - We relax the random response order assumption by investigating the performance of policies when response order is proportional to $\theta$. 
    For each patient, we compute $\bar{\theta}_{i} = \frac{1}{M} \sum_{j=1}^{M} \theta_{i,j}$, then let the order of the patients be proportional to $\bar{\theta}_{i}$. 
    That is, we construct $\sigma$ by selecting without replacement patients proportional to $\bar{\theta}_{i}$. 


    \item \textbf{Batched Response} - We explore the impact that control over the response order has through a batching framework. While we have assumed that $\sigma$ is random, in reality, healthcare administrators might be able to exercise control over $\sigma$, which can allow for partial control over the response order. We use two batches of patients, where order within a batch is randomized, and we defer a full description of the batching algorithm to Appendix~\ref{sec:batching}. 

    \item \textbf{Multi-Patient Capacity} - We relax the unit capacity assumption for policies by letting policies have a per-provider capacity, $c_{j}=2$. 
    We extend each policy to account for this; the modified policies are detailed in Appendix~\ref{sec:non_unit_policies}. 
    At a high level, each policy has a natural extension to the non-unit capacity scenario; for example, the pairwise policy requires a modification to the constraints of the linear program to account for this. 
\end{enumerate}

\begin{figure*}[ht!]
    \centering 
    \includegraphics[width=\textwidth]{figures/assumptions.pdf}
    \caption{We explore policy performance when relaxing three assumptions: we experiment with: policies that update online rather than being computed offline, response orders that are proportional to $\theta$ and are not uniformly random, response orders that are controllable by batching patients, and scenarios where providers have non-unit capacities. Across all three relaxations, we have the same ordering of policies regardless of the underlying assumptions. 
    We note that increasing the provider capacity results in homogenization because all policies increase match quality due to increased provider availability.}
    \label{fig:assumptions}
\end{figure*}
We evaluate the performance of policies under these relaxed assumptions when $\theta$ is uniformly distributed, and when $f$ is the uniform choice model for $p \in \{0.1,0.25,0.5,0.75,0.9\}$, while fixing $M=25$ and $N=100$. 
In Figure~\ref{fig:assumptions}, we find that relaxing assumptions minimally impact policy performance. 
Under all relaxations, the relative ordering of policies, given $p$, remains the same. 
Dynamic policies perform better than the pairwise policy, but despite their increased flexibility, are still worse than alternatives such as greedy and gradient descent. 
Similarly, even policies that can control the order $\sigma$ fail to improve performance compared to the gradient descent policy. 
Finally, increasing capacity leads to improved performance for the greedy policy, which occurs because provider scarcity is reduced and is equivalent to reducing the ratio of patients to providers (we demonstrate the latter claim through additional experiments in Appendix~\ref{sec:varied_capacity}). 
In addition to these experiments, in Appendix~\ref{sec:misspecification}, we demonstrate that our results are robust to misspecification in $\theta$ and $p$. 
\section{Analysis in a Semi-Synthetic Dataset}
\label{sec:real_world}
We analyze policy performance in a semi-synthetic dataset and show that the gradient descent policy can improve match quality by 13\%. 
We then discuss tradeoffs between various metrics, including match quality, match rate, and fairness. 


\subsection{Constructing a Semi-Synthetic Simulation}
\label{sec:semi_synthetic_data}
Beyond synthetic data, we work with a large healthcare provider in Connecticut to develop a simulation that reflects their system dynamics, leveraging real-world Medicare data and clinically grounded model parameters. 
We similarly average policy performance over $T$ rounds with randomized $\sigma$, but select the parameters $N$, $M$, $\theta$, and $f$ to accurately reflect their reality: 
\begin{enumerate}
    \item \textbf{$N$ and $M$} - We estimate $N=1225$ and $M=700$ from the average panel size and provider count of the organization.
    \item \textbf{Choice Model $f$} - We select the choice model based on previous work which shows that patients are low-effort decision makers who primarily make decisions based on geographic proximity~\citep{choosing_doctor,distance_rural}.
    We incorporate these two factors and let $f$ be a threshold model, with $\alpha$ representing the maximum distance patients are willing to travel. 
    \item \textbf{Match Quality Matrix $\theta$} - We construct $\theta$ to reflect  geographic proximity and the presence of comorbidities. 
    Geographic proximity is the top factor impacting patient match quality~\citep{washington_distance,distance_rural}, while patients with comorbidities are best served by providers with the corresponding specialized training. 
    Formally, for a patient $i$ and provider $j$, let $d_{i,j}$ be the distance between the patient and provider, and let $\beta_{i,j}$ denote whether patient $i$'s comorbidity and $j$'s specialty match.  
    For example, if patient $i$ has a heart-related comorbidity and provider $j$ has a cardiology specialty, then $\beta_{i,j}=1$. 
    The match quality is $\theta_{i,j} = \alpha + (1-\alpha)(\Delta \beta_{i,j} + (1-\Delta)(\frac{\bar{d}}{d_{i,j}}-1)$, where $\Delta$ and $\alpha$ are hyperparameters that modulate the impact of distance and comorbidities. 
    Here, $\Delta$ weighs between comorbidities and distance, and $\alpha$ sets the threshold match quality for a patient at distance $d_{i,j} = \bar{d}$.
    We provide details on our selection of hyperparameters in Appendix~\ref{sec:semi_synthetic_details}. 
\end{enumerate}

\subsection{Dataset Analysis}
\begin{figure*}[ht!]
    \centering 
    \includegraphics[width=\textwidth]{figures/patient_distro.pdf}
    \caption{We compare the distribution of patients (left), providers (middle), and the ratio of providers to patients (right). We find disparities between the distribution of patients and providers, revealing inherent difficulties in designing matching algorithms.} 
    \label{fig:distribution_patients}
\end{figure*}
To better understand our semi-synthetic dataset, we analyze the distribution of providers and patients. 
In Figure~\ref{fig:distribution_patients}, we show that the distribution of providers is more clustered in large cities, such as Hartford and New Haven, compared with the distribution of patients, matching results from prior work~\citep{rural_hospital_closure}. 
Moreover, plotting the ratio of providers to patients shows that most zip codes lack any providers while a few near metropolitan areas have a high density of providers (potentially corresponding to healthcare facilities). 
The distribution of patients and providers reveals inherent tensions for algorithm design; maximizing the match quality might lead to prioritization of ``easy-to-match'' patients near cities at the expense of more rural patients.  

\subsection{Match Rate and Match Quality Analysis}
\label{sec:semi_synthetic_results}
\begin{figure*}[ht!]
    \centering 
    \subfloat[\centering Comparison]{\includegraphics[width=0.45\textwidth]{figures/semi_synthetic_comparison.pdf}}
    \subfloat[\centering Comorbidities]{\includegraphics[width=0.45\textwidth]{figures/comorbidity_comparison.pdf}} \\ 
    \subfloat[\centering Match Rate vs. Population]{\includegraphics[width=0.45\textwidth]{figures/scatter_ct.pdf}}
    \subfloat[\centering Match Rate by Zipcode]{\includegraphics[width=0.45\textwidth]{figures/ct_map.pdf}}
    \caption{(a) We construct a semi-synthetic scenario with $N=1225$ and $M=700$. We find that the gradient descent policy performs best because of the large patient/provider ratio. 
    (b) The gradient descent policy achieves a higher match quality for patients with comorbidities because these patients match with specialist providers. (c) Geographically, the gradient descent policy tends to match patients from more populated zip codes, (d) with the highest match rates occurring in cities such as Hartford and New Haven.} 
    \label{fig:semi_synthetic}
\end{figure*}

In Figure~\ref{fig:semi_synthetic}a, we show that on the semi-synthetic dataset, the gradient descent policy achieves the highest match quality and outperforms alternatives by 13\%. 
This mirrors Figure~\ref{fig:phase}, where gradient descent is the best policy for large patient/provider ratios.  
Next, we analyze \textit{which} patients achieve high-quality matches under the gradient descent policy. 
We find that patients with comorbidities achieve higher match quality ($p<10^{-17}$) because they tend to match with specialist providers (Figure~\ref{fig:semi_synthetic}b).
Additionally, patients from more populated zip codes achieve higher match rates (Figure~\ref{fig:semi_synthetic}c), due to the aforementioned lack of providers in rural areas~\citep{rural_health_care}. 
Pictorially, this corresponds to high match rates in urban areas such as Hartford and New Haven (Figure~\ref{fig:semi_synthetic}d).
Our results demonstrate that even if a policy improves performance, it can still result in inequities on \textit{who} gets matched, and so practitioners should be aware of this when selecting assortment policies. 

\subsection{Metrics Beyond Match Rate and Match Quality}
We next explore policy performance beyond match quality and match rate by introducing metrics such as fairness and regret. 

\subsubsection{Fairness}
\label{sec:fair}
Fairness is critical to patient-provider matching to ensure equitable outcomes. 
To better understand this, we first introduce several definitions of fairness, each inspired by prior work from the fair matching literature~\citep{fairness_optimization,fair_minimax}. 
\begin{itemize}
    \item \textbf{Minimum Match Quality} - One natural objective is to maximize the minimum match quality across all matched patients. Such an objective is natural when administrators aim to maintain a certain level of care across patients. 
    \item \textbf{Variance in Match Quality} - Another approach minimizes variance in match quality to ensure that the discrepancy in match quality between matched patients is small. 
    \item \textbf{Match Quality Range} - One final approach to match quality fairness is to minimize the absolute difference between the maximum and minimum match quality for matched patients. 
\end{itemize}
Each notion of fairness has different implications for the optimal policy. 
For example, optimizing for the minimum match quality might lead policies to only offer assortments to high match quality patient-provider pairs, whereas minimizing the variance in match quality might yield many low-quality matches.  

\begin{figure*}[h!]
    \centering 
    \includegraphics[width=\textwidth]{figures/fairness_comparison.pdf}   
    \caption{We compare the performance of policies of both match rate (top) and match quality (bottom) and three definitions of fairness: minimum (left), variance (middle), and range (right) of match quality.  We find that a) the three definitions of fairness are aligned, as policies tend to do well or poorly on all three, and b) there are tradeoffs between match rate, match quality, and fairness.}
    \label{fig:fairness}
\end{figure*}

To understand how policies trade-off between match rate, match quality, and fairness, we plot each definition of fairness against both match rate (top) and match quality (bottom) in Figure~\ref{fig:fairness}.
We find that the three definitions of fairness are well aligned; across all definitions, the pairwise policy performs best, as it maximizes the minimum match quality, and minimizes the range and variance of match qualities.
Similarly, the greedy and gradient descent policies perform poorest across all definitions of fairness, though gradient descent performs better than greedy for variance. 
We additionally find that policies that perform better at either match rate or match quality tend to perform poorer at these fairness metrics. 
Intuitively, increasing the match quality or match rate requires matching patients with less suitable (lower match quality) providers, which leads to less fair outcomes but improved overall performance. 
Our overall takeaway is that no one policy performs best across all metrics, so practitioners need to carefully weigh metrics to determine which is best.  

\subsubsection{Patient Regret}
The assortment-based matching process can lead to patient regret if their initial assortment differs from the assortment offered during matching time. 
While not the primary driver of patient satisfaction, regret can nonetheless impact patient happiness because initially offered providers might no longer be available. 
To better understand patient regret, we first define it as the difference in the maximum match quality between the original assortment and the offered assortment, where the latter removes providers selected by patients. 
Formally, we can write regret for patient $i$ as $\max \mathbf{X}_{i} \odot \theta_{i} - \max \mathbf{X}_{i} \odot \theta_{i} \odot \mathbf{y}_{\sigma^{-1}_{i}}$. 
Here, $\mathbf{y}_{\sigma^{-1}_{i}}$ is the set of available providers when patient $i$ makes a decision. 

\begin{figure*}[h!]
    \centering 
    \includegraphics[width=\textwidth]{figures/regret.pdf}   
    \caption{We compare the regret (left) and assortment size (right) across policies and find that tailoring policies such as through gradient descent or pairwise can reduce regret. Moreover, we find a correlation between assortment size and regret because smaller assortments can allow for less overlap between patients.}
    \label{fig:regret}
\end{figure*}

To understand how policies perform for regret, in Figure~\ref{fig:regret}, we plot both patient regret and the average assortment size across policies. 
We find that policies such as pairwise suffer no regret because they offer disjoint menus across patients, while policies such as greedy suffer from the most regret because they do not tailor menus. 
We also see a general trend between assortment size and regret; as assortment sizes get larger, patient regret tends to increase as well, though the trend is imperfect.
Taken together, we find that no policy dominates across all metrics.
Improving the match quality might sacrifice fairness or lead to geographical disparities, while policies that perform well for fairness might lead to poor match rates or match qualities in some scenarios. 
\section{Conclusion and Discussion}
\label{sec:discussion}
The patient-provider relationship is foundational for quality healthcare, yet provider turnover frequently leads to patients being without providers. 
Our work studies how to re-match patients in such a scenario by reformulating the problem through assortment optimization. 
Our study proposes algorithmic approaches for patient-provider matching systems and offers analytical and empirical insights to inform system design.
Moreover, these policies broadly apply to one-shot matching under random response order, an understudied phenomenon within the assortment literature.
We conclude by providing recommendations for deploying assortment-based systems into practice and discussing extensions of our work. 

\subsection{Managerial Insights}
We translate our takeaways into real-world implications: 
\begin{enumerate}
    \item  Larger assortments become increasingly useful as patients become more selective because patients are more likely to reject matches, making it important to give more options. 
    Moreover, assortment sizes should tailor to patient behavior; the choice model and match probability impacts the best policy. 
    \item Various metrics are in tension when constructing assortments. In many situations, we cannot simultaneously maximize match rate and match quality, and incorporating fairness only complicates the picture. 
    Ultimately, these models provide a slate of options for practitioners to consider when weighing the priority of various metrics. 
    \item Healthcare administrators should be aware of the underlying factors that dictate \textit{which} patients get matched. For example, providers are scarce in rural areas, so assortments might naturally lead to low match rates in rural areas. Algorithm designers should work with administrators during development to minimize these biases.   
\end{enumerate}

\subsection{Broader Impact and Extensions}
Our work proposes policies for assortment optimization in the one-shot random response setting, and future work can extend these policies to other models within patient-provider matching and domains. 
Within patient-provider matching, natural extensions include incorporating  other metrics, including provider workload and total response times. 
Other extensions include modeling a dynamic pool of providers, where providers might enter and exit throughout the matching process and incorporating side constraints for individual patients or providers. 

Outside of patient-provider matching, our policies apply to a diverse class of problems with offline assortments and online random responses. 
Such a setting is a feature of many two-sided markets, and we present two examples. 
The first is assortment allocation in food delivery platforms; these platforms offer trips that drivers pick up and respond to in some random order, and batch offering eases the central planner's logistical burden. 
A second example is the matching between students and courses, where universities could offer each student an assortment of courses, and between schools and children, where school choice might lead counties to offer each child multiple options~\citep{school_choice}. 
These examples highlight the broader relevance of our problem setting and the applicability of our policies across these domains. 
%\pagebreak
 
\bibliographystyle{ACM-Reference-Format}
\bibliography{references}

\pagebreak
\section*{Acknowledgements}
We thank Aakash Lahoti, Keegan Harris, Kevin Li, Ariana Tang, Fei Fang, and George Chen for their comments. We thank the reviewers at ML4H for their feedback on a preliminary version of this work. We thank Susan Barrett and Karen Hall for their input from a health system perspective. Co-author Raman is supported by an NSF GRFP Fellowship.
\appendix

\section{Misspecification Experiments}
\label{sec:misspecification}
\paragraph{Misspecification in $p$}
\begin{figure*}
    \centering 
    \includegraphics[width=\textwidth]{figures/misspecification.pdf}
    \caption{We evaluate the performance of policies when the observed match probability, $\hat{p}$, differs from the true match probability, $p$. We find that the observed match probability has little impact on the relative ordering of policies.} 
    \label{fig:misspecification}
\end{figure*}
Our experiments in Section~\ref{sec:empirical} consider situations with known values of $p$, but real-world scenarios frequently involve misspecified choice models. 
To model this, we run policies with some observed match probability $\hat{p}$ while letting $p$ be the match probability. 
This setup allows us to simulate real-world scenarios where match probability is unknown. 
We compare policies across $\hat{p} \in \{0.1,0.25,0.5,0.75,0.9\}$ and $p \in \{0.1,0.25,0.5,0.75,0.9\}$ with $M=25$ and $N=100$.
In Figure~\ref{fig:misspecification}, we demonstrate that policies perform similarly across $\hat{p}$ for fixed $p$. 
This trend holds because the random, greedy, and pairwise policies are unimpacted by $\hat{p}$, while other policies are not sensitive to the exact choice of the match probability. 
These results imply our policy's robustness to their exact knowledge of the choice model. 

\paragraph{Variance in $p$ and $\theta$}
\begin{figure*}[h!]
    \centering 
    \includegraphics[width=\textwidth]{figures/theta_p_shift.pdf}    
    \caption{We compare the performance of policies when either $\theta$ or $p$ is misspecified by some noise $\epsilon$. We find no noticeable differences between policy performances even as noise size is increased; that is, the relative ordering of policies remains constant. }
    \label{fig:dynamic_p_theta}
\end{figure*}
We consider the impact that noisy $p$ and $\theta$ have on the performance of policies. 
For each parameter, we modify it, so $p_{i} = p+\epsilon$ for some unknown noise $\epsilon$, and similarly, we let $\theta = \hat{\theta} + \epsilon$, where we only observe $\hat{\theta}$. 
We let $\epsilon = \mathcal{N}(0,s^{2})$ for $s \in \{0,0.05,0.1,0.2\}$, and similarly for $\epsilon$, we let $\epsilon = \mathcal{N}(0,s^{2})$ for $s \in \{0,0.05,0.1,0.2\}$. 
We again let $M=25$ and $N=100$. 

In Figure~\ref{fig:dynamic_p_theta}, we demonstrate that policies perform similarly across values of $s$. 
For both $p$ and $\theta$, we find that the policies that perform well under the noiseless scenario similarly perform well in the noisy scenario. 
For example, in all scenarios, we find that the gradient descent policy performs best, while the pairwise policy performs worst. 
This indicates that our results are not very sensitive to the exact knowledge of $p$ and $\theta$. 

\section{Policies with Non-Unit Capacities}
\label{sec:non_unit_policies}
We develop variants of each policy that work under non-unit capacity $c_{j}$. 
For greedy, we leave the policy as is, as it is not dependent on the unit capacity assumption. 
For both the pairwise and group-based assumption, we relax the initial linear programming solution so it has $\sum_{i=1}^{N} \pi^{P}(\theta)_{i,j} \leq c_{j}$, but proceeds as normal otherwise. 
Finally, for the gradient descent policy, we modify the computation of $h(\mathbf{X})_{i,j}$. 
Recall that $h(\mathbf{X})_{i,j}$ is an upper bound on the availability of provider $j$ for patient $i$. 
Originally, at position $t$, the availability probability is $(1-p'_{j})^{t-1}$, which represents the probability that none of the first $t-1$ patients select provider $j$. 
Here, we now modify this so that no more than $c_{j}-1$ patients select provider $j$ in the first $t-1$ patients. 
We compute this by approximating the number of patients who select provider $j$ through a normal approximation $\mathcal{N}(p'_{j}(t-1),p'_{j}(1-p'_{j})(t-1))$. 
Then we want to compute the probability that $\mathbf{Pr}[\mathcal{N}((t-1)p',(t-1)p'(1-p')) \leq c_{j}-1]$.
This is equivalent to $\Phi(\frac{c_{j}-1-(t-1)p'}{\sqrt{(t-1)p'(1-p')}})$, which we then lower bound as: 
\begin{equation}
\left\{
\begin{array}{ll}
\displaystyle
\frac{z}{z^2 + 1} \cdot \frac{1}{\sqrt{2\pi}} e^{-z^2/2} & \text{if } c_{j} - 1 \leq (t - 1)p' \\\\
(1 - p')^{t - 1} & \text{otherwise}
\end{array}
\right.
\end{equation}
Here, $z = \frac{c_{j}-1-(t-1)p'}{\sqrt{(t-1)p'(1-p')}}$, and we use a bound on $\Phi(z)$~\citep{normal_bounds}. 
\section{Varying Provider Capacity}
\label{sec:varied_capacity}
\begin{figure*}[h!]
    \centering 
    \includegraphics[width=\textwidth]{figures/capacity.pdf}
    \includegraphics[width=\textwidth]{figures/capacity_providers.pdf}    

    \caption{We vary the capacity for providers from $1$ to $4$, and find that increased provider capacity leads to homogenization of policy performance. We find that the same effect can be accomplished when expanding the number of providers as well; having 25 providers each with a capacity of 4 yields similar results to having 100 providers. }
    \label{fig:varied_capacity}
\end{figure*}
Following the results in Section~\ref{sec:assumption_relaxation}, we run experiments to understand the impact of provider capacity on performance. 
We vary provider capacity from 1 to 4 (so each provider can take on 1 to 4 patients) when fixing $N=100$ and $M=25$, and plot this in Figure~\ref{fig:varied_capacity}.  
We find that as the amount of provider capacity is increased, policy performances are increasingly homogeneous. 
As providers have a larger capacity, individual assortments matter less, as all patients have a higher chance of matching with their most preferred provider. 
Moreover, in Figure~\ref{fig:varied_capacity} bottom, we find that increasing provider capacity produces a similar result to increasing the total number of providers.
For example, having $M=25$ providers each with a capacity of 4 produces similar results to having $M=100$ providers. 

\section{Batching Algorithm}
\label{sec:batching}
We study how to construct batches so to maximize the match quality. 
Formally, suppose we construct $L$ batches, $b_{1}, b_{2},\ldots,b_{L}$, where $b_{k} \subseteq [N]$, and we make offers first to batch $b_{1}$. 
Each set of batches modifies the response order distribution so that $\sigma$ is now distributed according to $\sigma \sim S(b_{1},b_{2},\ldots,b_{L})$, where $S$ is some probability distribution function.
Throughout this section, let $\mathrm{MQ}(\mathbf{X},\theta,f,\sigma) = \mathbb{E}_{\sigma}[\frac{1}{N} \sum_{t=1}^{N} f_{\sigma_{t}}(\mathbf{X} \odot \mathbf{y}^{(t)}) \cdot \theta_{\sigma_{t}}]$. 
If we can find some optimal ordering $\sigma^{*} = \argmax_{\sigma} \mathrm{MQ}(\mathbf{X},\theta,f,\sigma)$, then our goal is to set $b_{1},b_{2},\ldots,b_{L}$ so $\sigma^{*}$ and $S(b_{1},b_{2},\ldots,b_{L})$ are similar. 

\paragraph{Finding the optimal ordering $\sigma^{*}$} 
To understand $\sigma^{*}$, we first introduce a policy $\pi^{A}(\theta,\sigma)$, which improves the pairwise policy with knowledge of the order $\sigma$. 
Recall that $v(\theta)_{i} = j$ if $\pi^{P}(\theta)_{i,j} = 1$, and let $v^{-1}_{j} = i$ if $v(\theta)_{i}=j$ and $\sigma^{-1}_{i} = t$ if $\sigma_{t} = i$. 
We can compute the order $\sigma^{*}$ that optimizes $\mathrm{MQ}(\pi^{A}(\theta),\theta)$.
$\pi^{A}$ adds to the pairwise assortment based on knowledge of the response order $\sigma$. 
For example, if $\sigma^{-1}_{i} = 1$ and $\sigma^{-1}_{i'} = 2'$, then we know that offering $v(\theta)_{i}$ to provider $i'$ can only improve the match quality because $v(\theta)_{i}$ will only be available if unselected by patient $i$. 
More generally, we add on pairs $(i',v(\theta)_{i})$ based on the order $\sigma$ when $\sigma^{-1}_{i} < \sigma^{-1}_{i'}$. 
Intuitively, if patient $i$ comes earlier than patient $i'$ in $\sigma$, then it's safe to offer $v(\theta)_{i}$ to patient $i'$ as well because patient $i$ will have the first opportunity to claim $v(\theta)_{i}$.
More formally, we include $(i',v(\theta)_{i})$ in the augmented assortment if: a) patient $i'$ prefers $v(\theta)_{i}$ at least as much as their original provider under $\pi^P$, i.e., $\theta_{i', v_i} \geq \theta_{i', v_{i'}}$, and b) $i$ comes before $i'$ in $\sigma$, i.e., $\sigma^{-1}i < \sigma^{-1}{i'}$.
We focus on such a policy because it represents a natural extension of the pairwise policy to incorporate the order $\sigma$. 
We next characterize $\sigma^{*}$ for the policy $\pi^{A}$ then use this to construct the batches $b_{k}$: 
\begin{restatable}{lemma2}{thmordering} 
    Let $G=(V,E)$ be a graph with $N$ nodes such that there is an edge from $i$ to $i'$ if $\theta_{i,v(\theta)_{i}} \leq \theta_{i,v_{i'}}$ or $v(\theta)_{i} = -1$ and $v_{i'} \neq -1$. 
    Let $\sigma^{*}$ be the optimal ordering for the policy $\pi^{A}$.
    Then, traversing the nodes defined by $\sigma^{*}$ is a reverse topological ordering on $G$. 
\end{restatable}
Intuitively, the graph $G$ captures who envies whom: if there's an edge from $i$ to $i'$, then $i$ would prefer the provider assigned to $i'$, suggesting $i'$ should come earlier in the response to minimize blocking.
We additionally let $v(\theta)_{i} = -1$ for unmatched patients. 
We demonstrate this through casework on pairs of nodes $i,i'$ in $\sigma^{*}$. 

\paragraph{Constructing batches $b$}
We next construct $b_{1},b_{2},\ldots,b_{L}$, so $\sigma \sim S(b_{1},b_{2},\ldots,b_{L})$ is close to the optimal ordering $\sigma^{*}$ for $\pi^{A}$. 
The idea is to partition $\sigma^{*}$ into $L$ batches; then differences only emerge within batches rather than between batches.
We note that while we can approximate $\sigma^{*}$ through topological ordering, computing the optimal $\sigma^{*}$ is not guaranteed (as there can be many topological orderings). 
Once we have an approximation of $\sigma^{*}$, we can use dynamic programming to find the set of batches $L$ that minimizes the number of edges within a batch. 
By minimizing the number of edges within a batch, we reduce stochasticity in $\sigma$, as edges between batches have a fixed order. 
Formally, we can guarantee that, given knowledge of $\sigma^{*}$, we can approximate this through batching for sufficiently large $L$ under a set of regularity conditions: 
\begin{restatable}{theorem2}{optimalordering}
    Let $\sigma^{*}$ be an optimal ordering.
    Let $G=(V,E)$ be a graph with $N$ nodes such that node $i$ is connected to $i'$ if $\theta_{i,v(\theta)_{i}} \leq \theta_{i,v_{i'}}$ or $v(\theta)_{i} = -1$ and $v_{i'} \neq -1$. 
    Suppose that there exists a partition of $\sigma^{*}$ into $L$ batches, $b_{1},b_{2},\ldots,b_{L}$, such that for any partition $b_{k}$, no $i, i' \in b_{k}$ have a common descendant in $G=(V,E)$. 
    Then $\pi^{A}$ achieves the same match quality under partition $g$  as under the optimal ordering: 
    \begin{equation}
        \mathbb{E}_{\sigma \sim S(b_{1},b_{2},\ldots,b_{L})}[\mathrm{MQ}(\pi^{A}(\theta),\theta,\sigma)] = \mathrm{MQ}[\pi^{A}(\theta),\theta,\sigma^{*}] 
    \end{equation}
\end{restatable}
We prove this by demonstrating that, if we can control for the descendants of each patient, then we can fully control the order $\sigma$, and let it be equal to the optimal order $\sigma^{*}$. 
Here, the common descendant condition ensures that the only source of uncertainty for the ordering $\sigma$ is the order of the batches themselves rather than any uncertainty for ordering \textit{within} a batch. 
Overall, the takeaway is that in some situations, controlling the batches $b$ is equivalent to controlling the full order $\sigma^{*}$. 

\section{Semi-Synthetic Experiment Details}
\label{sec:semi_synthetic_details}
We compute $d_{i,j}$ by sampling patient locations from Connecticut zip codes and obtaining provider locations from a Medicare dataset~\citep{Medicare}. 
We obtain $\beta_{i,j}$ from prior work on comorbidity rates ~\citep{age_comorbidity} and using information on provider specialties from a Medicare dataset~\citep{Medicare}. 
We let $p=0.75$ because most patients are low-effort, and $\Delta=\alpha=\frac{1}{2}$ to balance proximity and comorbidities. Finally, we set $\bar{d}=20.2$ since prior work demonstrates the average distance threshold for patients is 20.2 miles~\citep{washington_distance}. 

\section{Proofs}
\label{sec:proofs}

\textbf{Theorem~\ref{thm:one_m}} Let $f$ be the uniform choice model with probability $p$.
Suppose $M=1$, and let $u_{1},u_{2},\ldots,u_{N}$ be a permutation of $\{1,\ldots,N\}$ such that $\theta_{u_{1}} \geq \theta_{u_{2}} \cdots \theta_{u_{N}}$.
Let $s$ be defined as follows:
\begin{equation}
   s = \argmax_{s} (1-(1-p)^{s}) \frac{\sum_{i=1}^{s} \theta_{u_{i},1}}{s}
\end{equation}
Then the policy which maximizes match rate is $\pi(\theta) = \mathbf{X} = \mathbf{1}$ while the policy which maximizes match quality is $\pi(\theta) =  \mathbf{X}, \mathbf{X}_{u_{1},1} = \mathbf{X}_{u_{2},1} \cdots \mathbf{X}_{u_{s},1} = 1$, where $\mathbf{X}_{i,1}$ is 0 otherwise. 

\begin{proof}

Suppose that $s$ patients are offered the provider $j=1$: $\lVert \mathbf{X} \rVert_{1} = s$. 
In the uniform choice model, the probability that any of the $s$ patients select the single provider is $p$, so the probability of provider $j=1$ matching is $1-(1-p)^{s}$. 
Then the assortment that maximizes the match rate has $p=N$, and so $\mathbf{X}=\mathbf{1}$. 

Next, we find the optimal match quality assortment. 
By symmetry in the response order, the chance that each patient is selected is equal, and so the match quality is 
\begin{equation}
    (1-(1-p)^{s}) \frac{\sum_{i=1}^{N} \theta_{i,1} \mathbf{X}_{i,1}}{s}
\end{equation}
For fixed $s$, the optimal assortment selects the $s$ largest values of $\theta_{i,1}$. 
Next, we note that larger values of $s$ increase the match probability $(1-(1-p)^{s})$, but could decrease the average match quality for the selected patient, $\frac{\sum_{i=1}^{N} \theta_{i,1} \mathbf{X}_{i,1}}{s}$. 
To balance between the two factors, we can iterate through values of $s$ and compute $(1-(1-p)^{s})$ and $\frac{\sum_{i=1}^{N} \theta_{i,1} \mathbf{X}_{i,1}}{s}$. 
\end{proof}

\textbf{Theorem~\ref{thm:greedy_match}}
Let $f$ be the uniform choice model with match probability $p$. 
Then 
\begin{equation}
    \mathrm{MR}(\pi^{R}(\theta),\theta,f)  \geq \min(p,\frac{M}{N})
\end{equation}
Moreover, there exists no policy $\pi'$ such that $\mathrm{MR}(\pi'(\theta),\theta,f) > \mathrm{MR}(\pi^{R}(\theta),\theta,f)$

\begin{proof}
     
    For the greedy policy, we note that $\pi^{R}(\theta)_{i} = \mathbf{1}$ for all patients $i$. 
    Therefore, the only situation where a patient is unmatched under the uniform choice model is either a) there are no available providers, or b) they fail to match with probability $1-p$. 
    For the former situation, we note that this only occurs when all $M$ providers match, meaning that the match rate is $\frac{M}{N}$. 
    For the latter situation, we note that if all providers fail to match, then each patient has a $p$ chance of matching.
    Let $\alpha$ be the probability that all providers are taken. 
    Then our overall match rate is $\mathrm{MR}(\pi^{R}(\theta),\theta,f) = \alpha \frac{M}{N} + (1-\alpha) p \geq \min(p,\frac{M}{N})$. 
    Moreover, no policy can perform better because no policy can expand upon the set of assortments offered to patients. 
\end{proof}

\textbf{Theorem~\ref{thm:pairwise_match}}
Let $f$ be the uniform choice model with match probability $p$. 
Then 
\begin{equation}
    \mathrm{MR}(\pi^{P}(\theta),\theta,f)  = p \frac{\min(M,N)}{N}
\end{equation}

\begin{proof}
     
    In the pairwise policy, there are $\min(M,N)$ pairs of patients and providers that are matched up. 
    Each is successfully matched with probability $p$ in the uniform choice model, yielding the overall match probability to be $p \frac{\min(M,N)}{N}$
\end{proof}

\begin{lemma}
    \label{lem:augment}
    Let $\pi$ be a policy that augments the pairwise policy; $\pi(\theta)_{i,j} \geq \pi^{P}(\theta)_{i,j}$ for any $\theta$ for all $i,j$. 
    Let $G = (V,E)$ be a directed graph with $N$ nodes such that nodes $i$ and $i'$ are connected if $\pi(\theta)_{i,v(\theta)_{i'}} = 1$ . 
    Here, $v(\theta)_{i} = j$ if $\pi^{P}(\theta)_{i,j} = 1$.
    Then we have that
    \begin{equation}
        \mathrm{MR}(\pi,\theta,f) = \mathrm{MR}(\pi^{P},\theta,f) 
    \end{equation}
    if each component in $G$ is a complete digraph.  
\end{lemma}
\begin{proof}
% We will first prove the reverse direction; that if $G$ is a complete digraph, then $\pi$ and $\pi^{P}$ have the same match rate. 
Consider patient $i$ in a component that is a connected component of size $k$. 
% This means that $\sum_{j} \pi(\theta)_{i,j} = 1$. 
% We will prove the set of providers unmatched and offered to patient $i$ is non-empty; that is $\lVert \pi(\theta)_{\sigma_{t}} \odot \mathbf{y}^{(t)} \rVert_{1} > 0$ for any ordering $\sigma$ and $\mathbf{y}^{(t)}$ with $\sigma_{t}=i$. 
% $\lVert \pi(\theta)_{\sigma_{t}} \odot \mathbf{y}^{(t)} \rVert_{1} = k$ when $\mathbf{y}^{(t)} = \mathbf{1}$, and decreases by one each time a provider from $\pi(\theta)_{\sigma_{t}}$ is selected. 
For any $j$ with $\pi(\theta)_{i,j}=1$, there exists $k-1$ other $i'$ with $\pi(\theta)_{i',j}=1$ because all patients within the complete graph have provider $j$ on their assortment. 
Therefore, at most $k-1$ other patients can select providers from $\pi(\theta)_{i}$, and therefore $\lVert \pi(\theta)_{\sigma_{t}} \odot \mathbf{y}^{(t)} \rVert_{1}$ can decrease by at most $k-1$ before $\sigma_{t}=i$. 
% Therefore, $\lVert \pi(\theta)_{\sigma_{t}} \odot \mathbf{y}^{(t)} \rVert_{1} \geq k-(k-1) = 1 > 0$. 
Moreover, because $\min(M,N)$ patients have a non-empty assortment, in expectation, we have that $\min(M,N)p$ patients match.%, and so the match rate for the policy $\pi$ is $p \frac{\min(M,N)}{N}$, which is the same as the pairwise policy. 

% Next, we will prove the forward direction: if the match rates are equal between $\pi$ and $\pi^{P}$, then $G$ consists of components that are complete digraphs. 
% Therefore, the only way for $\pi$ to achieve the same match rate as $\pi^{P}$ is to offer $\min(M,N)$ patients a non-empty assortment. 
% Consider some node $u$ in the graph, corresponding to a patient. 
% Suppose that patient $i$ has an assortment of size $k$, indicating that there exist $k$ edges from $i$ to some node. 
% Suppose that there exists a node, $i'$ such that $(i',i)$ is an edge, but $(i,i')$ is not. 
% Then consider the ordering that places $i$ after its neighbors and $i'$. 
% In such an ordering, suppose $i'$ selects $v(\theta)_{i}$ and each neighbor of $i$, $u$, select $v_{u}$. 
% This results in neither $i$ nor its neighbors being available when $i$ must select a patient, leaving an empty assortment.
% Therefore, if patient $i$ has $v_{i'}$ on its assortment, then $i'$ must have $v(\theta)_{i}$ on its assortment for the policy $\pi$. 

% Next, we consider the scenario where there exists a node $w$ such that $w$ is a two-hop neighbor of $i$ but not an immediate neighbor of $i$. 
% Suppose $w$ is adjacent to $i'$, so that $i'$ is on $w$'s assortment. 
% Order the patients such that $w$ comes first, then $i$'s neighbors, then $i$. 
% Let $w$ select $i'$, let $i'$ select $i$, and let all of $i$'s other neighbors select themselves.
% This results in an empty assortment for $i$. 
% Consider the $\min(M,N)$ non-empty assortments; the restriction of $G$ to only these patients has no two-hop neighbors that are not also one-hop neighbors. 
% % In other words, every node in $G$ only has one-hop neighbors.  
% Therefore, in any component, all neighbors are connected, and so $G$ consists of components that are complete digraphs. 
\end{proof}


\textbf{Theorem~\ref{thm:grouping_match}}
Let $f$ be the uniform choice model. Then 
\begin{equation}
    \mathrm{MR}(\pi^{G}(\theta),\theta,f) = \mathrm{MR}(\pi^{P}(\theta),\theta,f)
\end{equation} 

\begin{proof}
    
    We prove that the group-based policy matches the assumptions of Lemma~\ref{lem:augment}, and therefore, has the same match rate as the pairwise policy. 
    We note that, within a group $\mathbf{q}$, all patients share the same assortment. 
    Moreover, this assortment is the union of assortments of all pairwise assortments for each patient. 
    Therefore, we can view each group as a complete graph among all $i$ with $q_{i}=1$, and so the induced graph $G$ consists of complete digraphs. 
\end{proof}


\textbf{Theorem~\ref{thm:gradient_match}}: Let $f$ be the uniform choice model with match probability $p$. Then for any $p<1$, there exists $\theta$ such that the unique maximizer of $\langle g(h(\mathbf{X})),\theta \rangle$, $\mathbf{X}^{*}$, has a match rate: 

\begin{equation}
    \mathrm{MR}(\mathbf{X}^{*},\theta,f) < \mathrm{MR}(\pi^{P}(\theta),\theta,f)
\end{equation}

\begin{proof}
     
    We will construct an example where the unique maximizer performs worse than the pairwise policy. 
    Consider $N=M=2$, and let $\theta_{1} = [1,0]$ and $\theta_{2} = [1,\epsilon]$. 
    Here, $\langle g(h(\pi^{P}(\theta))),\theta \rangle = 1+\epsilon$. 
    Note that any policy that does not offer all providers to some patients performs worse than $\pi^{P}(\theta)$. 
    Next, we will show that $\mathbf{X}_{1} = [1,0], \mathbf{X}_{2} = [1,1]$ outperforms $\pi^{P}(\theta)$. 
    Here, $h(\mathbf{X})_{2,2} = 1$ and $h(\mathbf{X})_{*,1} = 1-\frac{p}{2}$. 
    Therefore, the expected match quality is $2-p+\frac{p}{2}\epsilon$. 
    Setting $\epsilon < \frac{1-p}{1-p/2}$ (which can always be done when $p<1$) yields that $\mathbf{X}$ is better; moreover, setting $\mathbf{X}_{1,2} = 1$ can only lead to a worse match quality. 
    Under $\mathbf{X}$, the match rate then goes from $p$ to $p-\frac{p^2}{4}$, showing that $\mathrm{MR}(\mathbf{X},\theta,f) < \mathrm{MR}(\pi^{P}(\theta),\theta,f)$. 
\end{proof}

% \textbf{Theorem~\ref{thm:gradient_match_lower_bound}}
%     \textbf{Theorem~\ref{thm:gradient_match_lower_bound}}: Let $f$ be the uniform choice model with match probability $p$. 
%     Let $g^{*}(h^{*}(\mathbf{X}))_{i,j} = \mathrm{Pr}[\argmax \ \theta_{i} \odot \mathbf{y}^{(\sigma^{-1}_{i})} = j]$. 
%     Then there always exists some maximizer $\mathbf{X}^{*}$ of $p \cdot \langle g^{*}(h^{*}(\mathbf{X})),\theta \rangle$, such that: 
%     \begin{equation}
%     \mathrm{MR}(\mathbf{X}^{*},\theta,f) \geq \frac{1}{2} \mathrm{MR}(\pi^{P}(\theta),\theta,f) 
%     \end{equation}

% \begin{proof}
%      
%     We will first demonstrate that any maximizer satisfies two properties, then demonstrate how these properties imply the aforementioned match rate. 
%     First, we will demonstrate that we can transform any maximizer of $g^{*}(h(^{*}\mathbf{X}))$,we can find an alternate maximizer $\mathbf{X}^{*}$ such that every provider is offered to a patient; $\lVert \mathbf{X}_{*,j} \rVert_{1}\geq 1 \forall j$. 
%     Suppose $\mathbf{X}$ is a maximizer and $j$ is not offered to any provider. 
%     Then we note that $h^{*}(\mathbf{X})_{*,j} = 0$. 
%     Suppose we offer provider $j$ to some patient $i$; then $h(\mathbf{X})_{i,j} = 1$, as $j$ is not offered to any other provider. 
%     Moreover, $g^{*}(h^{*}(\mathbf{X}))_{i,j}$ can only increase from letting $\mathbf{X}_{i,j} = 1$. 
%     Moreover, because only $h^{*}(\mathbf{X})_{i,j} = 1$ across $i$, we have that $g^{*}(h^{*}(\mathbf{X}))_{i',j}$ is unimpacted for all $i'$. 
%     Finally, suppose that $g^{*}(h^{*}(\mathbf{X}))_{i,j'}$ decreases; then $j$ must be preferred to $j'$, and so $\theta_{i,j} \geq \theta_{i,j'}$. 
%     Taken together, this implies that $g^{*}(h^{*}(\mathbf{X}))$ can only increase from offering provider $j$ to patient $i$. 

%     Second, we will demonstrate that there exists a maximizer so the number of providers offered to patient $i$ is at most the number of patients whose assortment intersects with that of patient $i$. 
%     Suppose that patient $i$ is offered providers $j_{1},j_{2},\ldots,j_{k}$, such that $\theta_{i,j_{1}} \geq \theta_{i,j_{2}} \geq \cdots \geq \theta_{i,j_{k}} $.
%     Let $i_{1},i_{2},\ldots,i_{k'}$ be the set of patients whose assortment contains at least one of $j_{1},j_{2},\ldots,j_{k}$; that is, for all $i \in \{i_{1},i_{2},\ldots,i_{k'}\}$, $\sum_{j \in \{j_{1},j_{2},\ldots,j_{k}\}} X_{i,j} \geq 1$. 
%     Then either $k' \geq k$ or we can repeatedly remove the least preferred provider $j_{k}$ without impacting the overall match quality until $k' \geq k$. 
%     This occurs because at most the top $k'-1$ providers will be unavailable, and so in the worst case, the patient will be selecting provider $k'$. 
%     If $k' < k$, then any providers $k'+1,k'+2,\ldots,k$ will never be selected, and so can be removed without penalty. 
%     This similarly applies to any subset $r < k$ of these providers; for any such subset, there must exist at least $r$ patients offered by those providers. 
%     Otherwise, the worst-off provider $r$ would never be selected, and therefore, we could similarly shrink the list. 

%     We finally show that if the preceding two properties hold, then the maximizer $\mathbf{X}^{*}$ has the property that $\mathbf{X}^{*}_{i,j} \geq \pi^{P}(\theta')_{i,j}$ for some $\theta'$ for all $i,j$.
%     Suppose patient $i$ is offered $k$ providers.
%     Then we know that at least $k$ patients must be offered those $k$ providers, and each provider is offered to at least one patient.
%     Select any $k$ of these patients and these $k$ providers, and let the patients be $i_{1},i_{2},\ldots,i_{k}$ and similarly let the providers be $j_{1},j_{2},\ldots,j_{k}$. 
%     Construct a new match quality matrix $\theta'$, so that $\theta'_{i_{k},j_{k}} =1$; note that here $\mathbf{X}^{*}_{i,j} \geq \pi^{P}(\theta')_{i,j}$. 
%     Finally, consider some patient $i'$ that was offered $k'$ providers. 
%     Suppose of these some subset $r$ of the providers was in the previous $k$ that were paired up. 
%     Then we know that even after removing those $r$ providers, we would still have $k'-r$ patients that are offered these $k'-r$ providers.  
%     We can therefore iteratively pair up patients and providers and construct the match quality matrix $\theta'$ as before until either all providers are matched or all patients are matched. 
    
%     To complete the proof, consider any assortment $\mathbf{X}$ that is a superset of some output of the pairwise policy. 
%     Each patient $i$ can either match with $v(\theta)_{i}(\theta)$, does not match with any provider, or has an empty assortment because patient $i'$ matched with $v(\theta)_{i}$. 
%     Each patient can block at most one other patient, and therefore, the number of matches is at least half that of the pairwise policy. 
% \end{proof}

\textbf{Theorem~\ref{thm:greedy}}
Let $f$ be the uniform choice model with match probability $p$. 
For any $p$ and $\epsilon$, there exists a $\theta$ such that 
\begin{equation}
    \mathrm{MQ}(\pi^{R}(\theta),\theta,f)  \leq \epsilon \mathrm{MQ}(\pi^{*}(\theta),\theta,f)
\end{equation}

\begin{proof}
    
    We construct a problem instance where the greedy policy achieves an $\epsilon$ fraction of the optimal match quality. 
    Let $N=M$ and construct $\theta$ as follows: let $\theta_{1,1} = 1$, while $\theta_{i,1} = 2 \Delta$ for $i \neq 1$, where $\Delta \leq \frac{1}{2}$. Let $\theta_{i,j} = \Delta$ for all $i$ and $j \neq 1$. 

    Let $\pi$ be the policy so that $\pi(\theta)_{i} = \mathbf{e}_{i}$. 
    The expected match quality for this policy is $p(\Delta (N-1) + 1)$; for patients 2 to $N$, it achieves an expected match quality of $p(\Delta (N-1))$ in total, while for patient $1$ it achieves an expected match quality of $p$. 
    Because $\pi^{*}(\theta)$ is optimal, we have  
    \begin{equation}
        p(\Delta (N-1) + 1) = \mathrm{MQ}(\pi(\theta),\theta,f) \leq \mathrm{MQ}(\pi^{*}(\theta),\theta,f)
    \end{equation}

    Next, consider the greedy policy where $\pi^{R}(\theta) = \mathbf{1}$.
    All patients prefer the provider $1$, and because of this, all patients have an equal likelihood of matching with the provider $1$. 
    Each patient has a $\frac{1}{N}$ chance of matching with provider $1$, and therefore the expected match quality for provider $1$ is $\frac{1}{N} + \frac{N-1}{N} (2 \Delta)$. 
    For all the $N-1$ other providers, we receive a reward of $\Delta$ upon matching, and each patient matches with probability $p$; therefore, we receive a total match quality of at most $\frac{1}{N} + 2  \frac{N-1}{N} \Delta + p(N-1) \Delta$. 
    Taking the ratio of the greedy and optimal policies yields the following: 
    \begin{align}
        \small 
        \frac{\frac{1}{N} + \frac{2\Delta(N-1)}{N} + p(N-1) \Delta}{p (\Delta (N-1) + 1)} 
         \leq \frac{\frac{1}{N} + 2\Delta + p(N-1) \Delta}{p(\Delta (N-1) + 1)}  \\ \leq \frac{\frac{1}{N} + 2\Delta + p(N-1) \Delta}{p} 
         \leq \frac{1}{Np} + \frac{2\Delta}{p} + N \Delta 
    \end{align}
    Finally, letting $N=\frac{3}{\epsilon p}$ and $\Delta = \min(\frac{\epsilon}{3N},\frac{p\epsilon}{6})$ yields that the approximation ratio is $\epsilon$. 
\end{proof}

\textbf{Theorem~\ref{thm:lp}}
Let $f$ be the uniform choice model with match probability $p$. 
Then 
\begin{equation}
    \mathrm{MQ}(\pi^{P}(\theta),\theta,f) \geq p \mathrm{MQ}(\pi^{*}(\theta),\theta,f) 
\end{equation}

\begin{proof}
     
    The pairwise policy constructs an assortment so each patient is only offered one provider, and no provider is offered to more than one patient. 
    Under this scenario, each patient offered a non-zero assortment matches with probability $p$; that is, if $\pi^{P}(\theta)_{i,j}=1$, then the expected match quality is $p \theta_{i,j}$. 
    The match quality is then $\frac{p}{N} \sum_{i=1}^{N} \pi^{p}(\theta)_{i,j} \theta_{i,j}$

    Next, we upper bound the performance of the optimal policy.
    The optimal policy achieves a match quality of 
    $\mathbb{E}_{\sigma}[\frac{1}{N} \sum_{t=1}^{N} f_{\sigma_{t}}\left(\pi^{*}(\theta)_{\sigma_{t}} \odot \mathbf{y}^{(t)}\right)  \cdot \theta_{\sigma_{t}}]$
    For any order $\sigma$, we have that no provider can be selected by two patients, and each patient can select at most one provider. 
    That is, for any order, we can bound the match quality as 
    \begin{align}
        \mathrm{MQ}(\pi^{*}(\theta),\theta,f) \leq \max\limits_{\mathbf{X}, \sum_{i=1}^{N} X_{i,j} \leq 1, \sum_{j=1}^{M} X_{i,j} \leq 1} \frac{1}{N} \sum_{i=1}^{N} \sum_{j=1}^{M} X_{i,j} \theta_{i,j} \\ = \frac{1}{N} \sum_{i=1}^{N} \sum_{j=1}^{M} \pi^{P}(\theta)_{i,j} \theta_{i,j}
    \end{align}
    The second inequality comes from the definition of $\pi^{P}(\theta)$. 
    Therefore: 
    \begin{equation}
        p\mathrm{MQ}(\pi^{*}(\theta),\theta,f)\leq \frac{p}{N} \sum_{i=1}^{N} \pi^{P}(\theta)_{i,j} \theta_{i,j} =  \mathrm{MQ}(\pi^{P}(\theta),\theta,f)
    \end{equation}
\end{proof}

\thmgrouping* 
\begin{proof}
     
    We first compute the expected match quality for the pairwise policy. 
    Each of the $M$ matched patients has a $p$ chance of selecting their assigned provider, and $\theta_{i,j} \leq 1$; therefore, $\mathbb{E}_{\sigma,\theta}[\sum_{t=1}^{N}  (f_{\sigma_{t}}\left(\pi^{P}(\theta)_{\sigma_{t}} \odot \mathbf{y}^{(t)}\right) \cdot \theta_{\sigma_{t}})] \leq pM$.

    Next, we compute the expected match quality for the greedy policy. 
    First, note that because $\theta_{i,j}$ are distributed i.i.d, each provider has an equal chance of being the most preferred and available provider for patient $i$. 
    Under a uniform choice model, the match probability and match quality are independent; by symmetry for each provider $j$, $N/M$ patients aim to match with provider $j$ ,and so there is a $(1-(1-p)^{N/M})$ chance that provider $j$ is matched. 
    Next, we need to compute $\mathbb{E}[\theta_{i,j} | \text{j is top avail. for i}]$. 
    Fix $i$, then by symmetry, each $j$ has an equal chance of being the top provider. 
    Moreover, because $\theta_{i,j}$ is uniformly distributed, each provider has an equal selection probability for each timestep. 
    Therefore, $\mathbb{E}[\mathbf{y}^{(t)}_{j}] = \beta_{t}$ for some coefficients $\beta_{j}$. 
    Next, let $\mathbb{E}[(\theta_{i})_{(j)}]$ be the $j$th largest value among $\theta_{i,*}$. 
    Then
    \begin{equation}
        \mathbb{E}[\theta_{i,j} | \text{j is top avail. for i}] = \frac{\sum_{j=1}^{M} \beta_{t}(1-\beta_{t})^{j-1} \mathbb{E}[(\theta_{i})_{(j)}]}{\sum_{j=1}^{M} \beta_{t}(1-\beta_{t})^{j-1}}
    \end{equation}

    By Chebyshev's sum inequality, we have that 
    \begin{equation}
        \frac{\sum_{j=1}^{M} \beta_{t}(1-\beta_{t})^{j-1} \mathbb{E}[(\theta_{i})_{(j)}]}{\sum_{j=1}^{M} \beta_{t}(1-\beta_{t})^{j-1}} \geq \frac{1}{M} \sum_{j=1}^{M} \mathbb{E}[(\theta_{i})_{(j)}] = \frac{1}{2}
    \end{equation}
    Summing across all providers gives the appropriate ratio.  
\end{proof}

\textbf{Theorem~\ref{thm:lower_bound}}
The following holds for any $\mathbf{X}$ when $f$ is the uniform choice model with probability $p$:
{\small \begin{equation}
    p \cdot \langle g(h(\mathbf{X})), \theta \rangle \leq \mathrm{MQ}(\mathbf{X},\theta,f)
\end{equation}}

\begin{proof}
     
    We compute a matrix, $\hat{\mathbf{Y}}$, such that $Y_{i,j}$ corresponds to the probability that patient $i$ is matched with $j$. 
    Then $\mathrm{MQ}(\mathbf{X},\theta,f) = \hat{\mathbf{Y}} \cdot \theta$. We then need to prove that
    $p\cdot \langle g(h(\mathbf{X})),\theta \rangle - \langle \hat{\mathbf{Y}},\theta \rangle \leq 0$
    To prove this, first let $\ell_{u_{i,j}} = (p \cdot g(h(\mathbf{X}))_{i,j} - \hat{Y}_{i,j})$. 
    We will first bound $\sum_{j=1}^{M} \ell_{u_{i,j}}$ for a fixed $i$. 
    Let $\theta_{i,u_{i,1}} \geq \theta_{i,u_{i,2}} \ldots \theta_{i,u_{i,M}}$ for some coefficients $u$. 
    For any $m$, $\sum_{j=1}^{m} \ell_{u_{i,j}} = \sum_{j=1}^{m} p \cdot g(h(\mathbf{X}))_{i,u_{i,j}}-\hat{\mathbf{Y}}_{i,u_{i,j}}$. 
    We can next compute $\sum_{j=1}^{m} p \cdot g(h(\mathbf{X}))_{i,u_{i,j}}$. 
    Recall from Section~\ref{sec:lower_bound} that we can explicitly compute $g(h(\mathbf{X}))_{i,j}$ through the function $d(n) = \frac{1}{N} \sum_{k=1}^{N} (1-p \frac{n-1}{N-1})^{k-1}$. 
    Then by the of definition $g(h(\mathbf{X})_{i,j})$, we have $p \sum_{j=1}^{m} g(h(\mathbf{X}))_{i,j} = p (\sum_{j=1}^{m} X_{i,u_{i,j}} d(\lVert \mathbf{X}_{*,u_{i,j}} \rVert_{1}) \prod_{l=1}^{j-1} (1-(\lVert \mathbf{X}_{*,u_{i,l}} \rVert_{1})))$. 
    Next, we note that $\sum_{j=1}^{m} \hat{\mathbf{Y}}_{i,u_{i,j}}$ represents the probability that any of the top-$m$ options are available and selected. 
    For any provider $j$, the probability that $j$ is available is at least $d(\lVert \mathbf{X}_{*,u_{i,j}} \rVert_{1})$ (see Section~\ref{sec:lower_bound}). 
    Therefore, 
    \begin{align}        
        \small 
        \sum_{j=1}^{m} \hat{\mathbf{Y}}_{i,u_{i,j}}  
        \geq p \mathrm{Pr}[\text{Provider 1 avail.} \vee \text{Provider 2 avail.} \cdots] \\ \geq p \sum_{j=1}^{m} X_{i,u_{i,j}} d(\lVert \mathbf{X}_{*,u_{i,j}} \rVert_{1}) \prod_{l=1}^{j-1} (1-(\lVert \mathbf{X}_{*,u_{i,j}} \rVert_{1})
    \end{align}
    Combining our statements for $\hat{\mathbf{Y}}$ and $p  \cdot g(h(\mathbf{X}))$ gives that $\sum_{j=1}^{m} X_{i,u_{i,j}} \hat{\mathbf{Y}}_{i,u_{i,j}} \geq \sum_{j=1}^{m} d(\lVert \mathbf{X}_{*,u_{i,m}} \rVert_{1}) \prod_{l=1}^{m-1} (1-(\lVert \mathbf{X}_{*,u_{i,m}} \rVert_{1}) = \sum_{j=1}^{m} p \cdot g(h(\mathbf{X}))_{i,u_{i,j}}$. 
    Therefore, we have that $\sum_{j=1}^{m} \ell_{u_{i,j}} \leq 0$. 
    Next, we will show that $p \cdot \langle g(h(\mathbf{X}))_{i}, \theta_{i} \rangle -  \langle \hat{\mathbf{Y}}_{i},\theta_{i}\rangle = \sum_{j=1}^{M} \ell_{u_{i,j}} \theta_{i,j} \leq 0$ for all $i$. 
    We will prove this inductively; first note that $l_{i,u_{i,1}} \theta_{i,u_{i,1}} \leq 0$
    Next, assume that $\sum_{j=1}^{M-1} \ell_{u_{i,j}} \theta_{i,u_{i,j}}$, then: 
    
    \begin{align}
        \small 
        \sum_{j=1}^{m} \ell_{u_{i,j}} \theta_{i,u_{i,j}}  = \sum_{j=1}^{m-1} \ell_{u_{i,j}} \theta_{i,u_{i,j}} + \ell_{m} \theta_{i,u_{i,m}} 
        \\ \leq \sum_{j=1}^{m-1} \ell_{u_{i,j}} \theta_{i,u_{i,j}} + \ell_{m} \theta_{i,u_{i,m-1}}  
        \\ \leq (\sum_{j=1}^{m-1} \ell_{u_{i,j}}) (\sum_{j=1}^{m-1} \theta_{i,u_{i,j}} - \theta_{i,u_{i,m-1}}) \leq 0
    \end{align}
Therefore, $p \cdot \langle g(h(\mathbf{X}))_{i}, \theta_{i} \rangle \leq 0$ for all $i$, and so summing across all $i$ gives $p \cdot \langle g(h(\mathbf{X})) ,\theta \rangle \leq 0$
\end{proof}

% \thmtightlowerbound*
% \begin{proof}
%     Let $\mathbf{X} = \pi^{P}(\theta)$ for some given $\theta$. 
%     Then because $\pi^{P}(\theta)$ is computed by solving a 1-1 bipartite matching problem, we know that $\sum_{i=1}^{N} X_{i,j} \leq 1$ for all $j$ and $\sum_{j=1}^{M} X_{i,j} \leq 1$ for all $i$. 
%     Next, we note that $0 \leq \lVert \mathbf{X}_{*,j} \rVert \leq 1$, and moreover, if $\mathbf{X}_{i,j} = 1$, then $\lVert \mathbf{X}_{*,j} \rVert = 1$. 
%     Because $d(1) = 1$, we have that $h(\mathbf{X})_{i,j} = \mathbf{X}_{i,j}$. 
%     Finally, because each row $i$ of $h(\mathbf{X})_{i,j}$ has at most one $j$ with $\mathbf{X}_{i,j} = 1$, and because $\bar{\mathbf{X}}_{i,u_{k}} = h(\mathbf{X})_{i,u_{k}} p \prod_{k'=1}^{k-1} (1-h(\mathbf{X})_{i,u_{k}})$, we have that $g(h(\mathbf{X})) = h(\mathbf{X}) = p \mathbf{X}$. 
%     Therefore, $\langle g(h(\mathbf{X})) ,\theta \rangle$ is exactly p times the value of the bipartite match. 
%     Finally, 
%     \begin{equation}
%         \mathbb{E}_{\sigma}[\sum_{t=1}^{N} (f_{\sigma_{t}}\left(\mathbf{X}_{\sigma_{t}} \odot \mathbf{y}^{(t)}\right)  \cdot \theta_{\sigma_{t}})]  = p \sum_{i=1}^{N} \sum_{j=1}^{M} \pi^{P}(\theta)_{i,j} \theta_{i,j} = p (f(g(\mathbf{X})) \cdot \theta).
%     \end{equation}
%     Therefore, the expected match quality for the pairwise policy $\pi^{P}(\theta)$ is the same as $p \langle \mathbf{X},\theta \rangle$.
% \end{proof}

\thmmnlmatch*
\begin{proof}
     We first prove the match rate equality, then prove the corresponding claim for match quality. 
    For match rate, we note that there are $\sum_{i} \mathbbm{1}[v(\theta)_{i} \geq 0]$ different patients who have a non-empty assortment (that is, patient $i$ has an assortment with $v(\theta)_{i}$). 
    Each patient matches with probability $\frac{\exp(\theta_{i,v(\theta)_{i}})}{\exp(\theta_{i,v(\theta)_{i}}) + \exp(\gamma)}$. 
    Summing over all $v(\theta)_{i} \geq 0$, and dividing by $N$ gives the average match rate. 
    We can bound the match quality as 
    {\small \begin{align}
        \mathbb{E}_{\sigma}[\frac{1}{N} \sum_{t=1}^{N} f_{\sigma_{t}}\left(\pi^{*}(\theta)_{\sigma_{t}} \odot \mathbf{y}^{(t)}\right)  \cdot \theta_{\sigma_{t}}] \\ \leq \max\limits_{\mathbf{X}, \sum_{i=1}^{N} X_{i,j} \leq 1, \sum_{j=1}^{M} X_{i,j} \leq 1} \frac{1}{N} \sum_{i=1}^{N} \sum_{j=1}^{M} X_{i,j} \theta_{i,j} \\ = \frac{1}{N} \sum_{i=1}^{N} \sum_{j=1}^{M} \pi^{P}(\theta)_{i,j} \theta_{i,j}
    \end{align}}
    Next, we note that each edge in the optimal linear program $v(\theta)_{i}$ exists with probability $\frac{\exp(\theta_{i,v(\theta)_{i}})}{\exp(\theta_{i,v(\theta)_{i}} + \exp(\gamma))}$ if $v(\theta)_{i} \geq 0$, and $0$ otherwise. 
    Therefore, we get the final claim by summing over all edges and taking the average. 
\end{proof}

\thmordering*
\begin{proof}
We first note that $G$ is acyclic; if $\theta_{i,v(\theta)_{i}} \leq \theta_{i,v_{i'}}$ and $\theta_{i',v_{i'}} \leq \theta_{i',v(\theta)_{i}}$, then swapping $v(\theta)_{i}$ and $v_{i'}$ improves the pairwise policy $\pi^{P}(\theta)$. 
We note that $\pi^{P}(\theta)$ is the optimal solution for $\sum_{i=1}^{N} \sum_{j=1}^{M} \pi^{P}(\theta) \theta_{i,j}$ by definition, and so such a swap would violate this property, implying that $G$ is acyclic. 

Next, suppose that $\sigma^{*}$ is not a reverse topological ordering. 
Then there exists nodes $i$ and $i'$, so that $i$ comes before $i'$ in $\sigma^{*}$ and there exists an edge from $i$ to $i'$. 
Consider the following cases: 
\begin{enumerate}
    \item \textbf{Case 1:} There exists no node $w$, between $i'$ and $i$ in $\sigma$ so that $w$ and $i'$ share an edge (in either direction). 
    Then $i$ and $i'$ can be placed in sequence ($i$ comes immediately before $i'$) without impacting the expected reward. 
    This is because all nodes $w$ between $i$ and $i'$ have no impact on the preferences of $i'$ because $\theta_{i',v_{i'}} \geq \theta_{i',v_{w}}$. 
    Moreover, swapping $i$ and $i'$ can only increase the expected match quality; placing $i$ before or after $i'$ has no impact on the match quality for patient $i'$ (as $\theta_{i',v_{i'}} \geq \theta_{i',v(\theta)_{i}}$), but placing $i$ after $i'$ can increase the expected match quality if $v(\theta)_{i}$ goes unmatched. 
    \item \textbf{Case 2:} Suppose there exists a node $w$, such that there is an edge from $w$ to $i'$ Then the pair $(w,i')$ violates the reverse topological order, and so we can recurse on this. Note that this is a smaller length pair of nodes within the ordering $\sigma^{*}$. Similarly, if there is an edge from $i$ to $w$, then $(i,w)$ violates the reverse topological order ,and so we can recurse on this
    \item \textbf{Case 3:} Suppose that there exists a node $w$ such that there are edges from $i'$ to $w$ and $w$ to $i$. Then this creates a cycle, breaking the acyclic property mentioned earlier. 
    \item \textbf{Case 4:} Suppose that there exists a node $w$ such that there is an edge from $i'$ to $w$. There exists no path from $i'$ to $i$ due to the acyclic nature of the graph. 
    Additionally, $w$ does not have an edge over $i$ or any of the ancestors of $i$. 
    Move $i'$ and all of its descendants before $i$ in the graph; there are no edges from the descendants of $i'$ to $i$. 
    Therefore, this flips the order of $i$ and $i'$ without impacting any of the descendants; because only $i$ and $i'$ are impacted by this change (as the descendants of $i'$ are not impacted by $i$), such a change can only increase the expected match quality, as $i$ can potentially match with $v_{i'}$. 
    \item \textbf{Case 5:} Suppose there exists a node $w$ such that there exists an edge from $w$ to $i$. 
    Due to the acyclic nature of the graph, there are no edges from $i'$ (or any of its ancestors) to $w$ or any of the ancestors of $i$. 
    Now consider moving $i$ and its ancestors after $i'$. 
    Because there are no edges from $i'$ to the ancestors of $i$, the result of the assortment from the ancestors of $i$ has no impact on $i'$. 
    Therefore, while patient $i$ can now potentially match with provider $v_{i'}$, none of the ancestors of $i$ are negatively impacted, improving aggregate match quality.  
\end{enumerate}
Our cases cover every scenario for nodes between $i$ and $i'$, and in all cases, we can make a minor change to the order $\sigma^{*}$ to bring it closer to a reverse topological ordering while improving match quality or recurse upon a subset of the ordering. 
\end{proof}

\optimalordering*
\begin{proof}
    Let $\sigma^{*}$ be the optimal ordering and let $\sigma$ be any ordering from batches $b_{1},b_{2},\ldots,b_{L}$. 
    Then we will prove that for any node $i$, the order of descendants of $i$ will be the same in $\sigma^{*}$ and $\sigma$. 

    To prove this, consider a node $i$. 
    Let the descendants of node $i$ in the optimal ordering be $d_{1},d_{2},\ldots,d_{i}$. 
    We note that in $\sigma$, each $d_{i'}$ is in a separate partition by properties of $b_{k}$. 
    Additionally, because $b_{1},b_{2},\ldots,b_{L}$ is a partition of $\sigma^{*}$, the order of $d_{1},d_{2},\ldots,d_{i}$ is maintained. 
    Therefore, for any node, the same set of ancestors is maintained through the partition $b_{k}$. 

    Next, we will show that the descendants of patient $i$ dictate what patient $i$ selects in $\pi^{A}$. 
    Let $Z_{1}, Z_{2},\ldots,Z_{N}$ be a set of Bernoulli random variables, each of which is 1 with probability $p$.
    If $Z_{t}=1$ then patient $\sigma_{t}$ will select the highest available provider in their assortment, and if $Z_{t}=0$, then patient $\sigma_{t}$ will select no provider. 
    Additionally, let $V_{1}, V_{2}, \ldots, V_{N}$ be a set of random variables, where $Z_{t}$ denotes the provider selected by patient $\sigma_{t}$ ($V_{t} = 0$ if $Z_{t}=0$). 
    Note that $V_{1}$ is only a function of $Z_{1}$; if $Z_{1}=0$, then $V_{1}=0$, and otherwise, $V_{1} = \argmax_{j} \theta_{\sigma_{1},j}$. 
    Next, let $D_{\sigma_{t}} \subseteq [N]$ be the set of descendants for patient $\sigma_{t}$; that is, $d \in D_{\sigma_{t}}$ means there exists a path from $\sigma_{t}$ to $d$ in $G$. 
    Suppose that $V_{\sigma_{t'}} = f(\{Z_{d}\}_{d \in D_{\sigma_{t'}}})$ for all $t' \leq t$. 
    Then patient $\sigma_{t+1}$ will only select some $j$ so $\theta_{\sigma_{t+1},j} \geq \theta_{\sigma_{t+1},v_{\sigma_{t+1}}}$. 
    By the construction of the graph $G$, this corresponds to edges connected to patient $\sigma_{t+1}$, and so $V_{t+1} = g(\{V_{d}\}_{d \in D_{\sigma_{t+1}})}$.
    Next, note that $V_{d} = f(\{Z_{d'}\}_{d' \in D_{d}})$. 
    Therefore, we can rewrite $V_{t+1}$ as 
    \begin{equation}
        V_{t+1} = g(\{f(\{Z_{d'}\}_{d' \in D_{d}})\}_{d \in D_{t+1})}.
    \end{equation}
    We finally collect like terms and note that the set of descendants of the children of $\sigma_{t+1}$ is the set of descendants of $\sigma_{t+1}$. 
    Therefore, we rewrite $V_{t+1}$ as 
    \begin{equation}
        V_{t+1} = g(\{Z_{d}\}_{d \in D_{t+1}})
    \end{equation}
    Therefore, $V_{t+1}$ only depends on the descendants of $t+1$.
    Therefore, if the ordering of descendants is fixed, and $Z_{1},Z_{2},\ldots,Z_{N}$ are decided a priori, then $V_{\sigma_{t}}$ is also fixed. 
    The ordering of descendants is the same between $\sigma \sim S(b_{1},b_{2},\ldots,b_{L})$ and $\sigma^{*}$, so both achieve the same match quality. 
\end{proof}

\end{document}
